%% ============================================================================
%% Chapter 14: Rate Limiting, Throttling & Quota Systems
%% Mastering APIs: From Zero to Production Guru
%% ============================================================================
\chapter{Rate Limiting, Throttling \& Quota Systems}
\label{ch:rate_limiting}

%% ----------------------------------------------------------------------------
%% Executive Summary
%% ----------------------------------------------------------------------------
Every API has a capacity, whether or not anyone has measured it. A service
that cannot say \emph{no} does not have unlimited generosity; it has an
unpriced contract with catastrophe, and the invoice arrives as a latency
spike, a connection-pool exhaustion, or a 3~a.m.\ page. Rate limiting is the
discipline of saying no \emph{early, cheaply, predictably, and
informatively} --- before the expensive parts of your stack have done work
they will have to abandon. This chapter treats it as what it is: a control
system with mathematics, a distributed-systems problem with race conditions,
and a client-facing contract with wire semantics. All three must be right,
and the third is the one most often forgotten.

We begin by separating five ideas that casual usage conflates:
\emph{rate limiting} (short-window admission control per identity),
\emph{throttling} (shaping the flow, usually by delaying rather than
rejecting), \emph{quotas} (long-window budgets that bill or ration usage
over hours, days, or billing periods), \emph{concurrency limits} (how many
requests may be \emph{in flight}, not how many may start), and \emph{load
shedding} (the system protecting \emph{itself}, dropping work by priority
when capacity is exceeded regardless of who sent it). Each answers a
different question, and mature platforms run all five simultaneously.

The algorithmic core is derived, not just stated. We prove that the fixed
window counter admits up to $2N$ requests inside a window-length interval
(Theorem~\ref{thm:ch14-fixed-window-2x}) and that this bound is tight; we
derive the sliding window counter's weighted estimator and bound its error
strictly by the previous window's count
(Proposition~\ref{prop:ch14-swc-error}); we prove the token bucket's
burst-plus-rate envelope $A(\Delta t) \le B + r\,\Delta t$
(Proposition~\ref{prop:ch14-tb-envelope}); and we derive GCRA --- the Generic
Cell Rate Algorithm that powered ATM networks and now powers
\texttt{nginx}, Envoy, and most serious gateways --- from its theoretical
arrival time update rule, proving its exact equivalence to the token bucket
(Proposition~\ref{prop:ch14-gcra-equiv}). Every claim is backed by an
executable test: the 2x boundary burst is not an anecdote here, it is a
passing pytest.

Distribution changes everything and nothing. The algorithms survive
replication; the \emph{state} does not. We dissect the read-modify-write
race that makes naive multi-replica limiters over-admit, show why a Redis
sorted set driven by a single Lua script is atomic by construction, and map
the cell-based architectures (one global cell, regional cells, local
limiters with gossip) onto their consistency--latency--blast-radius
trade-offs, including the fail-open versus fail-closed decision you must
make \emph{before} Redis dies, not during. Client-facing semantics get equal
rigor: 429 with \texttt{Retry-After} in both its legal forms, the IETF
\texttt{RateLimit-*} header fields draft, 503 for shedding, and quota
endpoints --- because a limit a developer cannot observe is a limit they
will route around~\cite{ietf_ratelimit_headers,rfc9110}. We close with
multi-dimensional limiting: per-key, per-IP, per-endpoint, cost-weighted
budgets, adaptive AIMD limits, and priority-aware shedding, all set in the
Northwind Robotics fleet-telemetry universe.

\begin{keyconceptbox}
A rate limiter is a \textbf{control system}, not a gate. Its setpoint is
capacity, its sensor is measured load, its actuator is admission, and its
contract with clients is the headers it emits. Choose the algorithm by its
burst and memory envelope (proved, not vibes), keep the shared-state update
atomic (one Lua script, one cell), and never reject silently: a 429 without
\texttt{Retry-After} is an outage with extra steps.
\end{keyconceptbox}

%% ----------------------------------------------------------------------------
\section{Learning Objectives \& Prerequisites}
\label{sec:ch14-objectives}

\subsection*{Learning Objectives}
After completing this chapter, its lab, and its exercises, you will be able
to:

\begin{enumerate}[label=\textbf{LO-\arabic*.}, leftmargin=2.6cm]
  \item Distinguish rate limiting, throttling, quotas, concurrency limits,
        and load shedding by the question each answers, the time window each
        governs, and the failure each prevents; identify which of the five a
        given production incident actually required.
  \item Derive, not just describe, the five classic admission algorithms ---
        fixed window, sliding window log, sliding window counter, token
        bucket, and GCRA --- including the fixed-window 2x boundary burst
        (Theorem~\ref{thm:ch14-fixed-window-2x}), the sliding-counter error
        bound (Proposition~\ref{prop:ch14-swc-error}), the token-bucket
        envelope (Proposition~\ref{prop:ch14-tb-envelope}), and the GCRA
        TAT update rule with its token-bucket equivalence
        (Proposition~\ref{prop:ch14-gcra-equiv}).
  \item Implement all five algorithms in typed Python~3.11 against an
        injected virtual clock, and prove their documented behaviors with
        deterministic, offline unit tests.
  \item Explain why naive read-modify-write limiting races across replicas,
        and implement an atomic distributed sliding-window limiter on Redis
        sorted sets driven by a single Lua script, verified offline against
        fakeredis under a thread storm.
  \item Design the failure policy for a limiter outage: fail-open,
        fail-closed, or degraded-local, selected per endpoint class from an
        explicit decision matrix rather than from adrenaline.
  \item Emit correct client-facing semantics: 429 with \texttt{Retry-After}
        (delay-seconds or HTTP-date), the IETF \texttt{RateLimit-Limit},
        \texttt{RateLimit-Remaining}, \texttt{RateLimit-Reset}, and
        \texttt{RateLimit-Policy} fields, 503 for load shedding, and quota
        introspection endpoints~\cite{ietf_ratelimit_headers,rfc9110}.
  \item Build FastAPI middleware enforcing per-API-key limits with accurate
        headers across window boundaries and a tested fail-open/fail-closed
        switch~\cite{fastapidocs}.
  \item Compose multi-dimensional limits --- per-key, per-IP, per-endpoint,
        cost-weighted, and adaptive (AIMD) --- and prioritize critical over
        batch traffic under pressure.
\end{enumerate}

\subsection*{Prerequisites}
This chapter assumes Chapters~0--7: the HTTP request/response model and
header semantics (Chapter~\ref{ch:http_wire}), consuming APIs and reading
error responses in Python (Chapter~\ref{ch:consuming_python}), REST resource
modeling (Chapter~\ref{ch:rest_style}), JSON payloads
(Chapter~\ref{ch:serialization}), API-key authentication basics
(Chapter~\ref{ch:client_security}), building and testing FastAPI services
with \texttt{TestClient} (Chapter~\ref{ch:fastapi}), and structured error
responses (Chapter~\ref{ch:problem_details}). Comfort with Big-O notation
and elementary algebra (floor functions, geometric bounds) is assumed; every
derivation is self-contained. The GraphQL cross-reference in the interview
vault points forward to Chapter~\ref{ch:graphql} but is not required here.

\subsection*{Setup}
All listings run offline on Windows with Python~3.11+. One-time setup from
PowerShell in a scratch directory:

\begin{minted}{text}
python -m venv .venv
.\.venv\Scripts\Activate.ps1
pip install fakeredis fastapi httpx pytest
\end{minted}

fakeredis emulates Redis (including Lua script execution) entirely
in-process, so the ``distributed'' listings run with no server, no
containers, and no network; Section~\ref{sec:ch14-deps} notes what changes
when you point the same code at a real \texttt{redis>=5} deployment.
%% ----------------------------------------------------------------------------
\section{Deep Technical Mechanics \& Architectural Concepts}
\label{sec:ch14-mechanics}

\subsection{Why Limit at All: Four Distinct Failures, Five Distinct Tools}
\label{subsec:ch14-why}

Rate limiting is not one technique but a family, and most production
incidents involving ``rate limits'' are actually failures to distinguish its
members. Four pressures make admission control non-optional:

\begin{enumerate}
  \item \textbf{Capacity protection.} Every downstream --- connection pool,
        worker count, database, a partner API --- has a knee beyond which
        throughput \emph{falls} as offered load rises (retry storms
        amplify this). Rejecting cheaply at the edge preserves the
        expensive interior for requests that will complete.
  \item \textbf{Fairness.} Tenants share capacity. Without per-identity
        accounting, one runaway client becomes everyone's outage; the
        noisy-neighbor problem is a limiting problem.
  \item \textbf{Cost control.} Egress, LLM tokens, metered third-party
        calls, and report renders all have marginal cost. An unpriced API
        is a subsidy for whoever automates against it first.
  \item \textbf{Abuse prevention.} Credential stuffing, scraping, and
        inventory hoarding are economic attacks; raising the attacker's
        cost per attempt is the defense of first resort
        (Chapter~\ref{ch:api_security}).
\end{enumerate}

The five tools map onto five different questions:

\begin{definition}[Rate limit]
A per-identity admission rule over a \emph{short} window (seconds to
minutes): \emph{``how fast may this caller start requests right now?''}
Enforced statefully per key; rejection is routine, cheap, and expected to be
retried by well-behaved clients.
\end{definition}

\begin{definition}[Throttling]
\emph{Shaping} rather than rejecting: delaying, dequeuing, or degrading
admitted work so the system's \emph{processing} rate stays at target. A
leaky bucket \emph{as a queue} throttles; a 429 \emph{limits}. Throttling
costs latency; limiting costs a rejection.
\end{definition}

\begin{definition}[Quota]
A \emph{long-window} budget (hour, day, billing month) that rations
aggregate consumption: \emph{``how much may this customer use in total?''}
Quotas are business constructs --- they encode plans, tiers, and cost ---
whereas rate limits are protection constructs. A client can be inside its
monthly quota and still be rate-limited this second.
\end{definition}

\begin{definition}[Concurrency limit]
A bound on requests \emph{in flight}, not requests started per unit time.
Ten requests that each hold a thread for 30 seconds stress a service more
than a hundred 5-millisecond requests; only a concurrency limiter (or bulkhead,
Chapter~\ref{ch:microservices}) sees the difference. Little's law connects
them: concurrency $=$ throughput $\times$ latency.
\end{definition}

\begin{definition}[Load shedding]
Self-protection independent of identity: when the system is saturated,
decline work --- by priority class, by cheapness to serve, or at random ---
to keep the admitted majority healthy. Shedding answers \emph{``which work
do we drop when we cannot do it all?''}; the honest response is 503 with
\texttt{Retry-After}~\cite{rfc9110}.
\end{definition}

\begin{warningbox}
Teams routinely deploy \emph{one} of the five and believe they have the
other four. A per-key rate limit does nothing against a legitimate,
in-limit client whose requests each take 45~seconds (need: concurrency
limit); a monthly quota does nothing against a 30-second burst that melts
the connection pool (need: rate limit); neither protects you from a
thundering herd of anonymous traffic (need: shedding). The audit question
is never ``do we have rate limiting?'' but ``which of the five questions
can this system currently answer?''
\end{warningbox}

\subsection{Fixed Window Counter --- and Its Provable 2x Flaw}
\label{subsec:ch14-fixed-window}

The fixed window counter is the limiter everyone writes first: partition
time into windows of length $W$ anchored at $t = 0$, keep one integer
counter per key per window index $\lfloor t/W \rfloor$, and admit while the
counter is below the limit $N$. It is $O(1)$ memory, trivially shardable,
and its \texttt{RateLimit-Reset} semantics are crisp (the window boundary).
It is also wrong at the edges, and the wrongness is a theorem, not a
folklore warning.

\begin{theorem}[Fixed-window boundary burst]
\label{thm:ch14-fixed-window-2x}
A fixed-window limiter configured for $N$ requests per window $W$ can admit
$2N$ requests inside a single real-time interval of length at most $W$ ---
indeed, inside an interval of length $2\varepsilon$ for any
$\varepsilon \in (0, W/2]$ --- and $2N$ is tight.
\end{theorem}

\begin{proof}
Let $kW$ be any window boundary. Fire $N$ requests at time
$t_0 = kW - \varepsilon$: all fall in window $k-1$, counter reaches $N$,
all admitted. Fire $N$ more at $t_1 = kW + \varepsilon$: window index has
rolled to $k$, the counter reset to zero, and all $N$ are admitted. The
$2N$ admissions all lie in $[t_0, t_1]$, an interval of length
$2\varepsilon \le W$. For tightness: any interval of length $W$ intersects
at most two fixed windows, and each contributes at most $N$ admissions, so
no such interval can contain more than $2N$.
\end{proof}

The practical reading: a client --- or a synchronized fleet of mobile
clients retrying after a connectivity blip, which is how this fires in real
life --- can double its nominal rate \emph{every} window, forever, by
straddling boundaries. Worse, the burst can be concentrated in an
arbitrarily small $\varepsilon$, so the instantaneous rate is unbounded:
$2N$ requests can arrive within milliseconds of each other. If your limit
exists to protect a connection pool or a partner API with a per-second
cap, the fixed window does not enforce the thing you think it enforces.
Listing~\ref{lst:ch14-tests} proves Theorem~\ref{thm:ch14-fixed-window-2x}
as an executable test (20 admissions in a 2-millisecond interval under a
``10 per second'' limit).

\subsection{Sliding Window Log: Exactness at $O(N)$ Memory}
\label{subsec:ch14-sliding-log}

The sliding window log removes the boundary entirely: store the timestamp
of every admitted request for the key, and on each arrival at time $t$
evict timestamps $\le t - W$; admit if fewer than $N$ remain. The window
now \emph{slides with the caller}, so no boundary exists to straddle, and
the invariant is exact:

\begin{proposition}[Sliding-log exactness]
\label{prop:ch14-swl-exact}
The sliding window log admits a request at $t$ if and only if strictly
fewer than $N$ previously admitted requests lie in $(t-W, t]$. Consequently
no interval of length $W$ ever contains more than $N$ admissions: the
$(i{+}N)$-th admission is always strictly more than $W$ after the $i$-th.
\end{proposition}

The cost is memory: $O(N)$ per key, per window. At $N = 10{,}000$ per
minute across a million API keys, the log is the state. This is why the
exact algorithm is rare at platform scale and why the next two algorithms
exist: they trade exactness for $O(1)$ state. A second, subtler cost is
that \texttt{Retry-After} becomes the time until the \emph{oldest} in-window
request expires --- perfectly computable, but it means your reset hint
varies with the client's own history rather than aligning to a wall-clock
boundary (Section~\ref{subsec:ch14-headers}).

\subsection{Sliding Window Counter: The Weighted Approximation}
\label{subsec:ch14-sliding-counter}

The sliding window counter keeps only two integers per key: the current
fixed window's count $c$ and the previous window's count $p$. At time $t$,
a fraction $f = (t - \lfloor t/W \rfloor W)/W$ of the current window has
elapsed, so a fraction $1-f$ of the previous window still overlaps the
true sliding window $[t-W, t]$. \emph{Assuming} the previous window's $p$
requests were spread uniformly across it, the overlapping portion is
estimated as $p(1-f)$, giving the estimator
\begin{equation}
  \hat{c}(t) \;=\; c(t) \;+\; p\,(1-f).
  \label{eq:ch14-swc}
\end{equation}
Admit when $\hat{c} < N$. This is the algorithm Cloudflare famously
analyzed at scale, finding sub-0.01\% disagreement with the exact log on
real traffic --- but their traffic is smooth, and the estimator's quality
is a bound, not a vibe:

\begin{proposition}[Sliding-counter error bound]
\label{prop:ch14-swc-error}
Let $o$ be the true number of previous-window requests inside $[t-W, t]$.
Then $|\hat{c} - c_{\mathrm{true}}| = |p(1-f) - o| \le p \cdot \max(f, 1-f)
< p \le N$.
\end{proposition}

\begin{proof}
$0 \le o \le p$ by definition, and the overlap region has length $(1-f)W$.
The error $|\,p(1-f) - o\,|$ is maximized at the extremes of $o$: if all
$p$ requests landed in the non-overlapping part of the previous window
($o=0$), the estimator \emph{overcounts} by $p(1-f)$; if all $p$ landed in
the overlap ($o=p$), it \emph{undercounts} by $p - p(1-f) = pf$. The larger
of the two is $p\max(f, 1-f) < p$, and $p \le N$ because the limiter was
enforcing during the previous window.
\end{proof}

Proposition~\ref{prop:ch14-swc-error} tells you exactly when to distrust
this algorithm: error scales with $p$, so the estimate is worst precisely
after a window that was saturated --- i.e., precisely when limiting
mattered. The adversarial case (all of last window's traffic bunched at its
end, as a retry fleet produces) makes the counter \emph{over-admit} early in
the next window; Listing~\ref{lst:ch14-tests} exhibits exactly that
over-admission as a test.

\subsection{Token Bucket: Burst Capacity as a Conserved Quantity}
\label{subsec:ch14-token-bucket}

The token bucket reparameterizes the problem from \emph{counting} to
\emph{budgeting}. A bucket holds at most $B$ tokens and refills
continuously at $r$ tokens per second; a request that finds $\ge 1$ token
consumes one and is admitted, else rejected. State per key is two scalars
(token level, last-refill timestamp) and the update is lazy --- computed on
arrival, no background timers:
\begin{equation}
  \mathrm{tokens}(t) \;=\; \min\!\bigl(B,\; \mathrm{tokens}(t_0) + r\,(t - t_0)\bigr).
  \label{eq:ch14-tb-refill}
\end{equation}

\begin{proposition}[Token-bucket envelope]
\label{prop:ch14-tb-envelope}
Over any interval $[t_0, t_1]$, the number of admitted requests satisfies
$A \le B + r\,(t_1 - t_0)$, and the bound is tight. Hence the long-run
admitted rate is at most $r$, while instantaneous bursts up to $B$ are
fully admitted.
\end{proposition}

\begin{proof}
Tokens are conserved: every admission removes one token, and the only
source of tokens is refill at rate $r$ into a bucket that started the
interval with at most $B$. Tokens available in the interval $\le B +
r(t_1 - t_0)$, bounding admissions; the bound is attained by arriving
exactly as tokens accumulate. Dividing by $\Delta t = t_1 - t_0$ gives
average rate $\le r + B/\Delta t \to r$ as $\Delta t \to \infty$.
\end{proof}

Proposition~\ref{prop:ch14-tb-envelope} is why the token bucket is the
default answer in interviews and in gateways: one parameter pair $(r, B)$
independently controls the \emph{sustained} rate and the \emph{burst}
tolerance, and the envelope is the tightest possible for any limiter with
those two properties (any limiter that admits a burst of $B$ and sustains
$r$ must admit at least this envelope). Figure~\ref{fig:ch14-token-bucket}
shows the mechanics.

\begin{figure}[ht]
\centering
\begin{tikzpicture}[font=\small, >=stealth]
  % ---- left: bucket schematic ----
  \draw[thick] (0,0) -- (0,2.6) (2.2,0) -- (2.2,2.6) (0,0) -- (2.2,0);
  \fill[blue!18] (0.04,0.04) rectangle (2.16,1.35);
  \foreach \x/\y in {0.35/0.3, 0.85/0.3, 1.35/0.3, 1.85/0.3,
                     0.6/0.75, 1.1/0.75, 1.6/0.75, 0.85/1.15, 1.35/1.15}{
    \draw[fill=darkblue!70] (\x,\y) circle (0.14);
  }
  \draw[decorate, decoration={brace, amplitude=5pt}, thick]
    (2.45,0) -- (2.45,2.6) node[midway, right=7pt, align=left]
    {capacity $B$\\ \footnotesize (burst tolerance)};
  \draw[->, thick, darkblue] (1.1,3.6) -- (1.1,2.7)
    node[midway, right, align=left] {refill: $r$ tokens/s};
  \node[align=center] at (1.1,-0.7) {\footnotesize token bucket};
  \draw[->, thick] (-2.6,1.3) -- (-0.15,1.3)
    node[midway, above, align=center] {requests};
  \draw[->, thick, packtorange] (2.35,0.5) -- (3.9,0.5)
    node[right, align=left] {admit: $-1$ token};
  \draw[->, thick, packtorange, dashed] (-0.15,2.2) -- (-1.3,2.9)
    node[left, align=right] {reject: 429};
  % ---- right: token level timeline ----
  \begin{scope}[shift={(6.4,0)}]
    \draw[->] (0,0) -- (5.6,0) node[below left] {\footnotesize $t$ (s)};
    \draw[->] (0,0) -- (0,3.0) node[above] {\footnotesize tokens};
    \draw[dashed, gray] (0,2.4) -- (5.4,2.4) node[pos=0, left] {\footnotesize $B$};
    % idle full -> burst drains -> refill at r -> cap at B -> trickle
    \draw[very thick, darkblue] (0,2.4) -- (1.0,2.4);
    \draw[very thick, packtorange] (1.0,2.4) -- (1.0,0.0);
    \draw[very thick, darkblue] (1.0,0.0) -- (2.8,2.4);
    \draw[very thick, darkblue] (2.8,2.4) -- (3.6,2.4);
    \draw[very thick, darkblue] (3.6,2.4) -- (3.6,2.0);
    \draw[very thick, darkblue] (3.6,2.0) -- (4.1,2.4);
    \draw[very thick, darkblue] (4.1,2.4) -- (5.4,2.4);
    \node[note] at (1.0,-0.35) {\footnotesize burst: $B$ at once};
    \node[note] at (2.2,0.9) {\footnotesize refill slope $r$};
    \node[note] at (4.9,2.05) {\footnotesize drip-fed admits};
    \foreach \t in {1.0, 3.6}{
      \draw[gray!60, dotted] (\t,0) -- (\t,2.4);
    }
  \end{scope}
\end{tikzpicture}
\caption{Token bucket mechanics. Left: refill at $r$ tokens/s into a bucket
capped at $B$; admission consumes one token, empty buckets reject. Right:
token level over time --- a full bucket absorbs a burst of $B$ instantly
(vertical drop), refills linearly at $r$, then admits only as fast as
tokens drip in. The area discipline of
Proposition~\ref{prop:ch14-tb-envelope} is visible: burst height $B$,
sustained slope $r$.}
\label{fig:ch14-token-bucket}
\end{figure}

\subsection{Leaky Bucket: Meter versus Queue}
\label{subsec:ch14-leaky-bucket}

``Leaky bucket'' names two different machines, and conflating them in a
design review is a reliable way to build the wrong one. \textbf{As a
meter}, the leaky bucket is a bucket that \emph{drains} at a constant rate
and overflows (rejects) when inflow exceeds outflow --- this is precisely
GCRA, derived next, and it is what \texttt{nginx}'s \texttt{limit\_req}
implements. \textbf{As a queue}, it is a FIFO buffer of capacity $Q$:
requests enqueue, and a worker drains the queue at exactly $r$ per second.
The two have opposite burst postures. The meter admits bursts (up to its
tolerance) and enforces only the average; the queue \emph{destroys} bursts
entirely, presenting a perfectly smooth $r$-per-second stream downstream at
the cost of queueing delay up to $Q/r$ seconds --- and of tying up a
connection or queue slot per waiting request. Use the meter to protect a
\emph{rate}; use the queue to protect a \emph{downstream that panics at
jitter} (a legacy mainframe interface, a partner with hard pacing). The
queue variant is throttling in the strict sense of
Definition~14.2: it shapes rather than rejects, until the queue itself
fills.

\subsection{GCRA: One Timestamp to Rule Them All}
\label{subsec:ch14-gcra}

The Generic Cell Rate Algorithm, standardized for ATM traffic policing and
rediscovered by every gateway author since, reduces rate limiting to a
single timestamp per key. Parameterize by the \emph{emission interval}
$T = 1/r$ (the minimum spacing between conforming requests at the sustained
rate) and a \emph{burst tolerance} $\tau$. State is the \emph{theoretical
arrival time} $\mathrm{TAT}$: when the next request would arrive if the
client were perfectly paced. On an arrival at $t_a$:
\begin{equation}
  \text{reject if } t_a < \mathrm{TAT} - \tau;
  \qquad
  \text{else admit and } \mathrm{TAT} \leftarrow \max(t_a, \mathrm{TAT}) + T.
  \label{eq:ch14-gcra}
\end{equation}
Read it operationally. The condition says: a request may arrive up to $\tau$
\emph{early} relative to the schedule; arriving earlier than that means the
client has consumed more burst credit than exists. The update says: each
admission pushes the schedule later by one emission interval, but if the
client was idle ($t_a > \mathrm{TAT}$) the schedule first re-anchors to the
present --- idle time does not accrue credit beyond $\tau$, because
$\max(t_a, \mathrm{TAT})$ discards any schedule that has fully lapsed.
Rejection yields a ready-made \texttt{Retry-After}: $\mathrm{TAT} - \tau - t_a$.

\begin{proposition}[GCRA $\equiv$ token bucket]
\label{prop:ch14-gcra-equiv}
GCRA with emission interval $T$ and tolerance $\tau$ admits exactly the
same request streams as a token bucket with rate $r = 1/T$ and capacity
$B = 1 + \tau/T$.
\end{proposition}

\begin{proof}[Proof sketch]
Define the token level from GCRA state as
$L(t) = \min\bigl(1 + \tau/T,\; 1 + (t + \tau - \mathrm{TAT})/T\bigr)$.
The GCRA admission condition $t_a \ge \mathrm{TAT} - \tau$ is algebraically
$L(t_a) \ge 1$. On admission, $\mathrm{TAT} \leftarrow \max(t_a,
\mathrm{TAT}) + T$ decrements $L$ by exactly $1$ (the $\max$ clamps $L$ at
capacity, matching the bucket's $\min(B, \cdot)$). Between arrivals,
$\mathrm{TAT}$ is constant and $L$ grows at rate $1/T = r$, matching
Equation~\eqref{eq:ch14-tb-refill}. Same state dynamics, same decisions;
Listing~\ref{lst:ch14-tests} verifies the equivalence on a seeded
adversarial stream.
\end{proof}

Why keep both? The token bucket parameterizes as humans think
(\emph{rate} and \emph{burst size}); GCRA parameterizes as schedulers think
(\emph{spacing} and \emph{jitter tolerance}) and stores its state as one
timestamp that serializes trivially into a single Redis field --- no
two-field tuple, no lazy-refill ordering hazards. Most serious gateways
(nginx, Envoy, Kong's default) are GCRA underneath whatever their
configuration calls it.

\subsection{Choosing an Algorithm}
\label{subsec:ch14-choosing}

\begin{table}[ht]
  \centering
  \small
  \begin{tabularx}{\textwidth}{l l l l X}
    \toprule
    \textbf{Algorithm} & \textbf{Memory/key} & \textbf{Burst behavior} &
    \textbf{Boundary} & \textbf{Choose when} \\
    \midrule
    Fixed window & $O(1)$ & up to $2N$ at edges & hard, exploitable &
      Crisp reset headers; rough protection; short windows \\
    Sliding log & $O(N)$ & none beyond $N$ & none (exact) &
      Correctness-critical, modest $N$ (auth attempts, payments) \\
    Sliding counter & $O(1)$ & bounded by Prop.~\ref{prop:ch14-swc-error} &
      smoothed & Huge key space, smooth traffic, tolerance for $\pm p$ error \\
    Token bucket & $O(1)$ & exactly $B$, then $r$ & none &
      General default; independent rate/burst knobs \\
    GCRA & 1 timestamp & exactly $1 + \tau/T$ & none &
      Single-field state; gateway parity; precise pacing \\
    \bottomrule
  \end{tabularx}
  \caption{Admission algorithm comparison. All rows enforce the same
  nominal rate; they differ in memory, burst envelope, and boundary
  behavior. GCRA and token bucket are provably equivalent
  (Proposition~\ref{prop:ch14-gcra-equiv}); the choice between them is
  operational, not behavioral.}
  \label{tab:ch14-algorithms}
\end{table}

\begin{figure}[ht]
\centering
\begin{tikzpicture}[
  font=\footnotesize, >=stealth, node distance=0.5cm and 0.9cm,
  decision/.style={diamond, draw=darkblue, thick, align=center,
    inner sep=1.2pt, text width=3.1cm},
  term/.style={rectangle, draw, fill=gray!15, rounded corners, align=center,
    minimum width=3.0cm, minimum height=0.85cm},
  start/.style={rectangle, draw, rounded corners, align=center,
    minimum width=3.0cm, minimum height=0.85cm}]
  \node[start] (start) {Need admission control\\ for identity $k$};
  \node[decision, right=of start] (replicas) {Multiple replicas\\ share the limit?};
  \node[decision, right=of replicas] (exact) {Exactness required?\\ \tiny (auth, payments)};
  \node[term, above right=0.1cm and 0.9cm of exact] (log) {Sliding window log\\ \tiny (Redis ZSET + Lua)};
  \node[decision, below right=0.1cm and 0.9cm of exact] (bursts) {Must tolerate\\ legit bursts?};
  \node[term, right=of bursts] (bucket) {Token bucket / GCRA\\ \tiny $B$ sized to legit burst};
  \node[term, below=of bursts] (counter) {Sliding counter or\\ fixed window \tiny (short $W$)};
  \node[term, above=of replicas, yshift=0.35cm] (cell) {Redis cell architecture\\ \tiny Figure~\ref{fig:ch14-distributed}};
  \draw[->] (start) -- (replicas);
  \draw[->] (replicas) -- node[above] {no} (exact);
  \draw[->] (replicas) -- node[left, pos=0.35] {yes} (cell);
  \draw[->] (cell.east) -| node[pos=0.2, above] {\tiny then pick algorithm} (log.north);
  \draw[->] (exact) -- node[above] {yes} (log);
  \draw[->] (exact) -- node[above] {no} (bursts);
  \draw[->] (bursts) -- node[above] {yes} (bucket);
  \draw[->] (bursts) -- node[above] {no} (counter);
\end{tikzpicture}
\caption{Algorithm selection flow. The first question is topological (shared
state?), the second semantic (exactness?), the third behavioral (bursts?).
``No bursts tolerated'' with $O(1)$ memory means accepting the fixed or
sliding counter's window artifacts consciously.}
\label{fig:ch14-decision}
\end{figure}
\subsection{Distribution: The Race, the Cell, and the Failure Policy}
\label{subsec:ch14-distributed}

A rate limiter's algorithm survives replication; its \emph{state} does not.
The moment two replicas enforce one logical limit, the question stops being
``which algorithm?'' and becomes ``where does the counter live, and who can
see it when?''

\subsubsection{The read-modify-write race}
The naive distributed fixed window --- \texttt{GET} the counter, compare,
\texttt{INCR} --- over-admits under any concurrency at all. With limit $N$
and counter at $N-1$, two requests landing on two replicas both read $N-1$,
both conclude ``admit,'' and the counter ends at $N+1$; under a retry storm
the overshoot grows with concurrency, not with request rate. Even the
apparently atomic \texttt{INCR}-then-compare variant (increment first,
reject when over) is subtly broken: rejected requests still \emph{consume}
counter increments, so an attacker (or a desperate retrying client) can
poison a window for everyone including itself, and the counter can run to
millions while admitting nothing. The general principle is the one
Chapter~\ref{ch:idempotency} applied to writes and
\cite{kleppmann2017ddia} applies to all shared state: \emph{check and
mutate must be one indivisible step against the authoritative copy}.

The standard solution is Redis plus Lua: Redis executes each Lua script
atomically on its single command-processing thread, so a script that
evicts, counts, conditionally inserts, and sets expiry is a compare-and-set
of arbitrary complexity with no client-side round trips
(Listing~\ref{lst:ch14-redis}). The sorted-set sliding log is the usual
payload: timestamps as scores, a UUID per request as member, one
\texttt{ZREMRANGEBYSCORE}/\texttt{ZCARD}/\texttt{ZADD}/\texttt{PEXPIRE}
sequence per admission decision.

\subsubsection{Cell architectures}
A single global Redis (or Redis Cluster) holding all limiter state is a
\emph{global cell}: perfectly consistent, but every admission pays a network
round trip and the cell is a blast-radius amplifier --- when it dies, every
replica's limiter dies simultaneously (Section~\ref{subsec:ch14-failopen}).
The alternatives are points on a consistency--latency frontier:

\begin{itemize}
  \item \textbf{Global cell.} Strongest semantics, highest per-request
        latency (colocate the cell with the edge tier to keep it
        sub-millisecond), largest correlated-failure domain. Right answer
        for low-QPS, high-stakes limits: login attempts, payment attempts,
        password reset.
  \item \textbf{Regional cells.} One cell per region; a key's limit is
        enforced only within the region serving it. A roaming client gets up
        to $R \times$ its nominal limit across $R$ regions --- usually
        acceptable for abuse control, unacceptable for billing-grade quotas.
        Mitigate by pinning keys to home cells via consistent hashing.
  \item \textbf{Local in-process limiters with periodic sync.} Each replica
        enforces $\mathrm{limit}/R$ locally (or a gossiped/share-weighted
        fraction) and reconciles asynchronously. Lowest latency and zero
        hot-path dependency; over-admission is bounded by the number of
        replicas times the sync lag in units of the window. This is
        essentially how Cloudflare's edge limiting works: approximate at the
        edge, exact nowhere, correct \emph{enough} everywhere. The error
        budget for this design is Proposition~\ref{prop:ch14-swc-error}
        thinking applied to topology instead of time.
\end{itemize}

Consistent hashing ties the pieces together: route limiter state for key
$k$ to cell $h(k) \bmod C$ so that adding or losing a cell remaps only
$1/C$ of keys (their counters reset --- a bounded, temporary under-limit,
which fails \emph{safe} for the system and merely generous to $1/C$ of
clients). Figure~\ref{fig:ch14-distributed} shows the shape.

\begin{figure}[ht]
\centering
\begin{tikzpicture}[
  font=\footnotesize, >=stealth, node distance=0.55cm and 1.0cm,
  edge/.style={rectangle, draw, rounded corners, align=center,
    minimum width=2.5cm, minimum height=0.8cm},
  cell/.style={rectangle, draw, fill=darkblue!10, rounded corners,
    align=center, minimum width=2.9cm, minimum height=0.9cm},
  client/.style={rectangle, draw, fill=packtorange!15, rounded corners,
    align=center, minimum width=2.2cm, minimum height=0.7cm}]
  \node[client] (c1) {robot fleet\\ \tiny key nr-demo-key-001};
  \node[client, below=of c1] (c2) {partner app\\ \tiny key nr-demo-key-002};
  \node[edge, right=1.6cm of c1] (e1) {edge replica 1\\ \tiny in-process fallback};
  \node[edge, below=of e1] (e2) {edge replica 2\\ \tiny in-process fallback};
  \node[edge, below=of e2] (e3) {edge replica 3\\ \tiny in-process fallback};
  \node[rectangle, draw, dashed, align=center, right=1.7cm of e2,
        minimum width=3.0cm, minimum height=0.9cm] (router)
        {consistent-hash router\\ \tiny $h(k) \bmod C$ on limiter key};
  \node[cell, above right=0.15cm and 1.6cm of router] (cellA)
        {Redis cell A\\ \tiny Lua: ZSET sliding log};
  \node[cell, below right=0.15cm and 1.6cm of router] (cellB)
        {Redis cell B\\ \tiny Lua: ZSET sliding log};
  \draw[->] (c1) -- (e1);
  \draw[->] (c2) -- (e3);
  \draw[->] (e1.east) -- ++(0.5,0) |- (router.west);
  \draw[->] (e2) -- (router);
  \draw[->] (e3.east) -- ++(0.5,0) |- (router.west);
  \draw[->] (router.east) -- ++(0.4,0) |- (cellA.west)
        node[pos=0.75, above] {\tiny $h(k)$ even};
  \draw[->] (router.east) -- ++(0.4,0) |- (cellB.west)
        node[pos=0.75, below] {\tiny $h(k)$ odd};
  \node[note, text=packtgray, align=center] at ($(e2.south) + (0,-0.55)$)
    {cell down $\to$ policy switch: fail-open / fail-closed / degraded-local};
\end{tikzpicture}
\caption{Distributed limiter architecture. Edge replicas consult one shared
cell per key, chosen by consistent hashing so cell loss remaps only $1/C$
of keyspace; the Lua script makes check-and-admit atomic inside a cell. The
dashed policy annotation is mandatory, not decorative: every replica must
know what it does when the cell is unreachable
(Section~\ref{subsec:ch14-failopen}).}
\label{fig:ch14-distributed}
\end{figure}

\subsubsection{Failure modes: fail-open, fail-closed, or degrade}
\label{subsec:ch14-failopen}
A limiter whose state store is unreachable cannot count. What it does next
is a \emph{business decision wearing an engineering costume}, and it must be
made per endpoint class, in advance:

\begin{table}[ht]
  \centering
  \small
  \begin{tabularx}{\textwidth}{l X X X}
    \toprule
    \textbf{Policy} & \textbf{Behavior on limiter outage} &
    \textbf{Blast radius of the choice} & \textbf{Right for} \\
    \midrule
    Fail-open & Serve everything; emit no unverifiable headers &
      Abuse window until recovery; possible cascade if the outage
      \emph{was} load & Read-heavy public APIs, telemetry ingest \\
    Fail-closed & Reject (503 + \texttt{Retry-After}) &
      Total API outage for the duration; safe against abuse &
      Login, password reset, payment authorization \\
    Degraded-local & Fall back to per-replica limit $\mathrm{limit}/R$ &
      Bounded over-admission ($\le R\times$), no outage &
      The default for most multi-replica services \\
    \bottomrule
  \end{tabularx}
  \caption{Failure-policy decision matrix. Degraded-local dominates when a
  cheap local limiter is available; the fail-open/fail-closed binary is for
  endpoints where either abuse or availability is existentially worse.}
  \label{tab:ch14-failmatrix}
\end{table}

\begin{warningbox}
A fail-closed limiter with no exception is a single point of total failure:
when Redis dies, the API dies with it --- and because health checks usually
pass through the same middleware, your load balancer may agree that
everything is down and refuse to help. If you fail closed, exempt the
health endpoint, keep the failure \texttt{Retry-After} short, and alert on
limiter-unavailable as a severity-1 condition in its own right
\cite{nygard2018releaseit}.
\end{warningbox}

\subsection{Client-Facing Semantics: Making Limits Legible}
\label{subsec:ch14-headers}

A limit developers cannot observe will be treated as a bug --- theirs, then
yours. The wire contract has four moving parts.

\textbf{429 and Retry-After.} Rejection is status 429 \emph{Too Many
Requests}. RFC~9110 defines \texttt{Retry-After} in two legal forms: a
non-negative integer of \emph{delay-seconds}
(\texttt{Retry-After: 17}) or an absolute \emph{HTTP-date}
(\texttt{Retry-After: Thu, 03 Sep 2026 21:00:00 GMT})~\cite{rfc9110}.
Emit delay-seconds: HTTP-date costs you clock-sync reasoning on both sides
and buys nothing a client cannot compute from the integer. Fractional waits
must be ceiled --- a client that retries at $0.4$s when you meant $0.44$s
will be rejected again and will (correctly) conclude your header is noise.

\textbf{The IETF RateLimit fields.} The HTTPAPI working group's draft
standardizes four fields~\cite{ietf_ratelimit_headers}:
\texttt{RateLimit-Limit} (the quota for the window, e.g.\ \texttt{100}),
\texttt{RateLimit-Remaining} (budget left \emph{after} this response),
\texttt{RateLimit-Reset} (seconds until the window's quota fully resets ---
note: delay-seconds, \emph{not} a timestamp; several major providers shipped
epoch-seconds variants before the draft settled, so document yours), and
\texttt{RateLimit-Policy} (the rule itself, e.g.\ \texttt{"100;w=60"} for
100 per 60 seconds). Emit them on \emph{successes} too: a client that can
watch \texttt{Remaining} decay can self-throttle before ever seeing a 429,
which is the entire point. Headers you cannot compute truthfully (during a
fail-open episode) must be omitted, never faked --- Listing~\ref{lst:ch14-middleware}
implements exactly this rule.

\textbf{503 for shedding.} When the system sheds load for its own
protection (not because \emph{this client} exceeded \emph{its} limit), the
honest status is 503 \emph{Service Unavailable} with
\texttt{Retry-After}~\cite{rfc9110}. Conflating shedding with limiting
(``429 for everyone during the incident'') misleads well-behaved clients
into believing they did something wrong and misroutes their backoff logic:
a 429 says \emph{slow down, you specifically}; a 503 says \emph{come back
later, everyone}.

\textbf{Quota endpoints.} Long-window budgets deserve introspection:
\texttt{GET /v1/usage} returning the period's consumption, cap, and reset
timestamp turns a billing surprise into a dashboard. Stripe-style
idempotent, header-rich, self-describing limiting is what separates a
platform from an obstacle course.

\begin{examplebox}
Northwind Robotics publishes \texttt{RateLimit-Policy: "600;w=60"} on its
telemetry ingest API. A fleet-management integrator watches
\texttt{RateLimit-Remaining} cross 60 and linearly spaces its remaining
batch across the seconds left in the window; during a regional Redis cell
outage the middleware fails open \emph{and omits} all four headers, which
the integrator's SDK reads as ``limiter degraded, apply local pacing.'' No
page, no ticket, no surprise invoice. That is what legible limits buy.
\end{examplebox}

\subsection{Multi-Dimensional Limiting}
\label{subsec:ch14-multidim}

Real systems limit along several axes at once, because no single identity
captures every abuse shape:

\begin{itemize}
  \item \textbf{Per-API-key:} the billing identity; the primary axis for
        fairness and plan enforcement.
  \item \textbf{Per-IP:} the pre-authentication axis (you must limit
        \emph{before} you can know the key --- otherwise login and key
        validation are free to attack). Behind NATs and carrier-grade NAT,
        one IP may be a ten-thousand-person enterprise (fairness disaster
        if the limit is tight) or one IP may be a botnet of thousands
        (useless if the limit is loose). Per-IP limits are a coarse outer
        net, never the only one --- the case study in
        Section~\ref{sec:ch14-case} is precisely this failure.
  \item \textbf{Per-endpoint:} a global request budget per route, defending
        the route's specific scarce resource regardless of which keys
        consume it.
  \item \textbf{Cost-based:} charge each request a \emph{weight}
        proportional to its true cost --- a bulk export might cost 10 units
        against a read's 1 --- so the budget is denominated in work
        (Listing~\ref{lst:ch14-cost}). GraphQL makes this mandatory rather
        than optional: the \emph{query}, not the route, determines cost, so
        limiters there score the parsed query's field weights and depth
        (Chapter~\ref{ch:graphql}).
  \item \textbf{Adaptive:} AIMD --- additive increase, multiplicative
        decrease, TCP's gift to every backpressure problem. Raise a client's
        limit additively while the service is healthy ($L \leftarrow L + a$),
        slash it multiplicatively on overload signals --- p99 latency
        crossing target, 5xx rate, queue depth --- ($L \leftarrow \beta L$,
        $\beta \approx 0.5$--$0.9$). AIMD converges to the largest limit the
        system can currently serve and backs off faster than it advances,
        which is exactly the asymmetry safety wants. It pairs naturally with
        adaptive \emph{concurrency} limits on the server side.
  \item \textbf{Priority-aware shedding:} under pressure, admit critical
        traffic (health, control plane, paying interactive reads) and shed
        batch/analytics first. Priority must come from an authenticated
        attribute of the request (key tier, route class), never from a
        client-supplied header --- a \texttt{X-Priority: high} you honor is
        a \texttt{X-Priority: high} your abusers send.
\end{itemize}

Expensive endpoints deserve explicit mention because they are where
dimensionless limits go to die: report generation, CSV/PDF exports, bulk
ingest, and search export endpoints can each cost $10$--$1000\times$ a
median request. Protect them with (a) a per-endpoint cost weight, (b) a
separate, tighter concurrency cap (two concurrent exports, queue the rest),
and (c) where possible, an asynchronous shape instead --- accept the job,
return 202 with a status URL, and let the \emph{job queue} apply its own
pacing (Chapter~\ref{ch:async_apis}).
%% ----------------------------------------------------------------------------
\section{Production-Ready Code Implementation}
\label{sec:ch14-code}

Everything below is typed Python~3.11, offline-first, and deterministic:
every limiter takes an injected \texttt{Clock}, every timeline is seeded,
and no test reads wall-clock time. The code is set in the Northwind
Robotics universe --- a fleet-telemetry platform whose API keys all begin
\texttt{nr-demo-} (synthetic; no real credentials appear anywhere in this
book). Save the listings as indicated and the printed commands reproduce
every claim in this chapter.

\subsection{Listing 14.1 --- Five Algorithms, One Contract}
\label{subsec:ch14-listing-limiters}

\texttt{rate\_limiters.py} implements all five admission algorithms behind a
single \texttt{check(at=None) -> Decision} contract, where
\texttt{Decision} carries the three things middleware will need: admit flag,
remaining budget, and a \texttt{Retry-After} hint. The \texttt{VirtualClock}
is the determinism backbone: tests advance time explicitly, so behavior is
identical on every machine, every run.

\noindent\textbf{Listing 14.1 --- \texttt{rate\_limiters.py}: reference
implementations with injected clock.}
\label{lst:ch14-limiters}
\begin{minted}{python}
"""Reference implementations of the five classic rate-limiting algorithms.

Every limiter takes an injected ``Clock`` so behavior is fully deterministic
under test: no wall-clock reads anywhere. Time is in float seconds.

Algorithms implemented:
  * FixedWindowLimiter      -- counter per floor(t / window) bucket
  * SlidingWindowLogLimiter -- exact timestamp log
  * SlidingWindowCounterLimiter -- weighted previous/current window estimate
  * TokenBucketLimiter      -- capacity B, refill rate r
  * GCRALimiter             -- generic cell rate algorithm (leaky bucket as
                               a meter), parameterized by emission interval T
                               and burst tolerance tau

All limiters expose the same ``allow(at=None) -> bool`` contract and are
single-process reference models; Listing 14.4 lifts the sliding log to Redis.
"""

from __future__ import annotations

import math
from collections import deque
from typing import Deque, Optional, Protocol


class Clock(Protocol):
    """Anything that reports the current time in float seconds."""

    def now(self) -> float: ...


class VirtualClock:
    """Deterministic, manually advanced clock for tests and simulations."""

    def __init__(self, start: float = 0.0) -> None:
        self._t = float(start)

    def now(self) -> float:
        return self._t

    def advance(self, seconds: float) -> None:
        if seconds < 0:
            raise ValueError("cannot move a clock backwards")
        self._t += seconds


class Decision:
    """Result of an admission check: admit flag plus header metadata."""

    def __init__(self, allowed: bool, remaining: int, retry_after: float) -> None:
        self.allowed = allowed
        self.remaining = remaining
        self.retry_after = max(0.0, retry_after)

    def __repr__(self) -> str:  # pragma: no cover - debugging aid
        return (
            f"Decision(allowed={self.allowed}, remaining={self.remaining}, "
            f"retry_after={self.retry_after:.3f})"
        )


class FixedWindowLimiter:
    """Allow at most ``limit`` requests per fixed window of ``window`` seconds.

    Memory: O(1). Flaw: up to ``2 * limit`` requests can be admitted inside
    any real-time interval of length ``window`` (see Theorem 14.1).
    """

    def __init__(self, limit: int, window: float, clock: Clock) -> None:
        if limit <= 0 or window <= 0:
            raise ValueError("limit and window must be positive")
        self.limit = limit
        self.window = window
        self.clock = clock
        self._window_index = -1
        self._count = 0

    def check(self, at: Optional[float] = None) -> Decision:
        t = self.clock.now() if at is None else at
        index = math.floor(t / self.window)
        if index != self._window_index:
            self._window_index = index
            self._count = 0
        window_end = (index + 1) * self.window
        if self._count < self.limit:
            self._count += 1
            return Decision(True, self.limit - self._count, 0.0)
        return Decision(False, 0, window_end - t)


class SlidingWindowLogLimiter:
    """Exact limiter: keep the timestamp of every admitted request.

    Memory: O(limit) per key. Never admits more than ``limit`` requests in
    any interval of length ``window`` -- the gold standard for correctness.
    """

    def __init__(self, limit: int, window: float, clock: Clock) -> None:
        if limit <= 0 or window <= 0:
            raise ValueError("limit and window must be positive")
        self.limit = limit
        self.window = window
        self.clock = clock
        self._log: Deque[float] = deque()

    def check(self, at: Optional[float] = None) -> Decision:
        t = self.clock.now() if at is None else at
        cutoff = t - self.window
        while self._log and self._log[0] <= cutoff:
            self._log.popleft()
        if len(self._log) < self.limit:
            self._log.append(t)
            return Decision(True, self.limit - len(self._log), 0.0)
        oldest = self._log[0]
        return Decision(False, 0, oldest + self.window - t)


class SlidingWindowCounterLimiter:
    """Weighted approximation of the sliding log with O(1) memory.

    Estimate for the true count in [t - window, t]:
        estimate = current + previous * (1 - elapsed_fraction)
    where elapsed_fraction is how far ``t`` has progressed into the current
    fixed window. The absolute error versus the exact log is strictly less
    than the previous window's count (Proposition 14.2).
    """

    def __init__(self, limit: int, window: float, clock: Clock) -> None:
        if limit <= 0 or window <= 0:
            raise ValueError("limit and window must be positive")
        self.limit = limit
        self.window = window
        self.clock = clock
        self._prev_index = -1
        self._prev_count = 0
        self._curr_index = -1
        self._curr_count = 0

    def check(self, at: Optional[float] = None) -> Decision:
        t = self.clock.now() if at is None else at
        index = math.floor(t / self.window)
        if index != self._curr_index:
            self._prev_index = self._curr_index
            self._prev_count = self._curr_count if self._curr_index == index - 1 else 0
            self._curr_index = index
            self._curr_count = 0
        elapsed_fraction = (t - index * self.window) / self.window
        previous = self._prev_count if self._prev_index == index - 1 else 0
        estimate = self._curr_count + previous * (1.0 - elapsed_fraction)
        window_end = (index + 1) * self.window
        if estimate < self.limit:
            self._curr_count += 1
            return Decision(True, max(0, int(self.limit - estimate - 1)), 0.0)
        return Decision(False, 0, window_end - t)


class TokenBucketLimiter:
    """Token bucket: capacity ``burst`` tokens, refilled at ``rate`` per second.

    Bursts up to ``burst`` are admitted instantly; the long-run admitted rate
    never exceeds ``rate`` (Proposition 14.3).
    """

    def __init__(self, rate: float, burst: float, clock: Clock) -> None:
        if rate <= 0 or burst <= 0:
            raise ValueError("rate and burst must be positive")
        self.rate = rate
        self.burst = burst
        self.clock = clock
        self._tokens = float(burst)
        self._last = clock.now()

    def _refill(self, t: float) -> None:
        if t < self._last:
            raise ValueError("timestamps must be non-decreasing")
        self._tokens = min(self.burst, self._tokens + (t - self._last) * self.rate)
        self._last = t

    def check(self, at: Optional[float] = None, cost: float = 1.0) -> Decision:
        if cost <= 0:
            raise ValueError("cost must be positive")
        t = self.clock.now() if at is None else at
        self._refill(t)
        if self._tokens >= cost:
            self._tokens -= cost
            return Decision(True, int(self._tokens), 0.0)
        deficit = cost - self._tokens
        return Decision(False, int(self._tokens), deficit / self.rate)


class GCRALimiter:
    """Generic Cell Rate Algorithm (leaky bucket used as a meter).

    Parameters: ``rate`` requests per second sustained, ``burst`` requests
    tolerated instantly. Emission interval T = 1 / rate; burst tolerance
    tau = (burst - 1) * T. State is a single timestamp: the theoretical
    arrival time (TAT) of the next conforming request.
    """

    def __init__(self, rate: float, burst: int, clock: Clock) -> None:
        if rate <= 0 or burst <= 0:
            raise ValueError("rate and burst must be positive")
        self.rate = float(rate)
        self.emission_interval = 1.0 / rate  # T
        self.tolerance = (burst - 1) * self.emission_interval  # tau
        self.clock = clock
        self._tat = float("-inf")

    def check(self, at: Optional[float] = None) -> Decision:
        t = self.clock.now() if at is None else at
        arrival_limit = self._tat - self.tolerance
        if t < arrival_limit:
            return Decision(False, 0, arrival_limit - t)
        self._tat = max(t, self._tat) + self.emission_interval
        # Remaining burst: admits still possible right now before the
        # not-before time t >= TAT - tau would be violated.
        headroom = t + self.tolerance - self._tat
        remaining = 1 + int(headroom // self.emission_interval) if headroom >= 0 else 0
        return Decision(True, remaining, 0.0)
\end{minted}

\subsection{Listing 14.2 --- The Theorems, as Executable Tests}
\label{subsec:ch14-listing-tests}

\texttt{test\_rate\_limiters.py} is the chapter's proof sheet:
Theorem~\ref{thm:ch14-fixed-window-2x} (20 admissions in 2~ms under a
``10/s'' limit), Proposition~\ref{prop:ch14-swl-exact} (the sliding-log
interval invariant), Proposition~\ref{prop:ch14-swc-error} (adversarial
bunching produces exactly the predicted over-admission),
Proposition~\ref{prop:ch14-tb-envelope} (the envelope never violated over
2000 steps), and Proposition~\ref{prop:ch14-gcra-equiv} (GCRA and token
bucket agree on a 500-event seeded adversarial stream).

\noindent\textbf{Listing 14.2 --- \texttt{test\_rate\_limiters.py}:
deterministic proofs of documented behavior.}
\label{lst:ch14-tests}
\begin{minted}{python}
"""Deterministic unit tests for the five reference limiters.

Every test injects a VirtualClock; nothing here reads wall time. The
fixed-window double-burst flaw and the GCRA/token-bucket equivalence are
proven as executable tests, mirroring the chapter's theorems.
"""

from __future__ import annotations

import pytest

from rate_limiters import (
    FixedWindowLimiter,
    GCRALimiter,
    SlidingWindowCounterLimiter,
    SlidingWindowLogLimiter,
    TokenBucketLimiter,
    VirtualClock,
)


def run_burst(limiter, n: int, at: float) -> int:
    """Fire n requests at the same instant; return how many were admitted."""
    return sum(1 for _ in range(n) if limiter.check(at=at).allowed)


class TestFixedWindow:
    def test_admits_up_to_limit_per_window(self) -> None:
        clock = VirtualClock()
        limiter = FixedWindowLimiter(limit=3, window=10.0, clock=clock)
        assert run_burst(limiter, 3, at=1.0) == 3
        assert not limiter.check(at=2.0).allowed
        assert not limiter.check(at=9.9).allowed
        # A fresh window restores the full quota.
        assert limiter.check(at=10.0).allowed

    def test_double_burst_at_boundary(self) -> None:
        """Executable proof of Theorem 14.1: a fixed-window limiter configured
        for N per W admits up to 2N requests inside a single real-time
        interval of length W that straddles a boundary."""
        clock = VirtualClock()
        limiter = FixedWindowLimiter(limit=10, window=1.0, clock=clock)
        admitted = 0
        # 10 requests at the very end of window 0 ...
        admitted += run_burst(limiter, 10, at=0.999)
        # ... and 10 more at the very start of window 1. Both windows admit
        # their full quota, yet all 20 requests fall inside [0.999, 1.001],
        # a real interval of 2 ms -- far shorter than the nominal window.
        admitted += run_burst(limiter, 10, at=1.001)
        assert admitted == 20  # 2x the nominal limit
        assert 1.001 - 0.999 < 1.0


class TestSlidingWindowLog:
    def test_exactness(self) -> None:
        clock = VirtualClock()
        limiter = SlidingWindowLogLimiter(limit=3, window=10.0, clock=clock)
        # Three requests mid-window saturate it; crossing the fixed
        # boundary at t=10 does NOT restore quota, because the window
        # slides with the caller.
        assert run_burst(limiter, 5, at=5.0) == 3
        assert run_burst(limiter, 5, at=10.0) == 0
        assert run_burst(limiter, 5, at=14.999) == 0
        # At t=15.0 the first request has fully aged out.
        assert limiter.check(at=15.0).allowed

    def test_no_interval_exceeds_limit(self) -> None:
        """For any request pattern, the admitted set never contains more
        than ``limit`` events inside any half-open interval of length
        ``window``: the (i + limit)-th admission is always strictly more
        than ``window`` after the i-th."""
        clock = VirtualClock()
        limiter = SlidingWindowLogLimiter(limit=4, window=2.0, clock=clock)
        admitted: list[float] = []
        t = 0.0
        for _ in range(400):
            t += 0.05
            if limiter.check(at=t).allowed:
                admitted.append(t)
        for i in range(len(admitted) - 4):
            assert admitted[i + 4] - admitted[i] > 2.0 - 1e-9

    def test_retry_after_points_at_oldest_expiry(self) -> None:
        clock = VirtualClock()
        limiter = SlidingWindowLogLimiter(limit=2, window=10.0, clock=clock)
        limiter.check(at=1.0)
        limiter.check(at=3.0)
        decision = limiter.check(at=4.0)
        assert not decision.allowed
        assert decision.retry_after == pytest.approx(7.0)  # 1.0 + 10.0 - 4.0


class TestSlidingWindowCounter:
    def test_matches_exact_when_previous_window_empty(self) -> None:
        clock = VirtualClock()
        limiter = SlidingWindowCounterLimiter(limit=3, window=10.0, clock=clock)
        assert run_burst(limiter, 5, at=15.0) == 3

    def test_weights_previous_window_by_elapsed_fraction(self) -> None:
        clock = VirtualClock()
        limiter = SlidingWindowCounterLimiter(limit=10, window=1.0, clock=clock)
        # Saturate window 0 entirely at its start.
        assert run_burst(limiter, 10, at=0.01) == 10
        # Halfway through window 1 the estimate is 0 + 10 * (1 - 0.5) = 5,
        # so exactly 5 more requests are admitted before t = 1.5.
        assert run_burst(limiter, 10, at=1.5) == 5

    def test_error_bound_holds(self) -> None:
        """Executable check of Proposition 14.2: the absolute error of the
        weighted estimate is strictly less than the previous window count."""
        clock = VirtualClock()
        counter = SlidingWindowCounterLimiter(limit=8, window=1.0, clock=clock)
        exact = SlidingWindowLogLimiter(limit=8, window=1.0, clock=clock)
        # Adversarial: all of window 0's traffic bunched at its very end.
        run_burst(counter, 8, at=0.99)
        run_burst(exact, 8, at=0.99)
        admitted_counter = run_burst(counter, 20, at=1.01)
        admitted_exact = run_burst(exact, 20, at=1.01)
        # The exact log admits 0 here (all 8 requests still inside the
        # sliding window); the counter over-admits 1 because it assumes
        # window-0 traffic was spread uniformly. Error 1 < previous
        # window count 8, as the bound predicts.
        assert admitted_exact == 0
        assert admitted_counter == 1
        assert abs(admitted_counter - admitted_exact) < 8


class TestTokenBucket:
    def test_burst_equals_capacity(self) -> None:
        clock = VirtualClock()
        limiter = TokenBucketLimiter(rate=10.0, burst=10, clock=clock)
        assert run_burst(limiter, 100, at=0.0) == 10
        # Half a second of silence refills exactly five tokens.
        clock.advance(0.5)
        assert run_burst(limiter, 100, at=clock.now()) == 5

    def test_long_run_rate_never_exceeds_refill(self) -> None:
        """Executable check of Proposition 14.3: over any interval [t0, t1]
        the admitted count is at most burst + rate * (t1 - t0)."""
        clock = VirtualClock()
        limiter = TokenBucketLimiter(rate=10.0, burst=10, clock=clock)
        admitted = 0
        t = 0.0
        for _ in range(2000):
            t += 0.02
            if limiter.check(at=t).allowed:
                admitted += 1
        assert admitted <= 10 + 10.0 * t + 1e-9

    def test_retry_after_is_deficit_over_rate(self) -> None:
        clock = VirtualClock()
        limiter = TokenBucketLimiter(rate=4.0, burst=2, clock=clock)
        run_burst(limiter, 2, at=0.0)
        decision = limiter.check(at=0.0)
        assert not decision.allowed
        assert decision.retry_after == pytest.approx(0.25)


class TestGCRA:
    def test_equivalent_to_token_bucket(self) -> None:
        """GCRA with emission interval T and tolerance tau = (B-1)T admits
        exactly the same streams as a token bucket with capacity B and rate
        1/T. Proven by construction in Section 14.3 (GCRA); checked here on a
        seeded adversarial stream."""
        import random

        rng = random.Random(14)
        gcra = GCRALimiter(rate=10.0, burst=10, clock=VirtualClock())
        bucket = TokenBucketLimiter(rate=10.0, burst=10, clock=VirtualClock())
        t = 0.0
        for _ in range(500):
            t += rng.choice([0.0, 0.005, 0.02, 0.2, 1.0])
            assert gcra.check(at=t).allowed == bucket.check(at=t).allowed

    def test_steady_stream_at_rate_always_conforms(self) -> None:
        clock = VirtualClock()
        limiter = GCRALimiter(rate=5.0, burst=3, clock=clock)
        for i in range(100):
            assert limiter.check(at=i * 0.2).allowed

    def test_retry_after_is_tat_minus_tolerance_minus_now(self) -> None:
        clock = VirtualClock()
        limiter = GCRALimiter(rate=10.0, burst=2, clock=clock)  # T=0.1, tau=0.1
        assert limiter.check(at=0.0).allowed   # TAT: 0.1
        assert limiter.check(at=0.0).allowed   # TAT: 0.2
        decision = limiter.check(at=0.0)       # arrival limit = 0.2 - 0.1
        assert not decision.allowed
        assert decision.retry_after == pytest.approx(0.1)
\end{minted}

Run it:

\begin{minted}{text}
pytest test_rate_limiters.py -q
# 14 passed
\end{minted}

\subsection{Listing 14.3 --- Simulation Harness: One Timeline, Five Verdicts}
\label{subsec:ch14-listing-sim}

Claims about burst behavior should be measured, not asserted.
\texttt{burst\_simulation.py} replays one seeded synthetic timeline --- a
25-request synchronized-retry burst straddling the $t{=}1.0$ fixed-window
boundary, followed by a steady 8~rps stream --- through all five limiters
configured at the same nominal 10~requests/second, and prints the
per-250~ms admission counts that produce Figure~\ref{fig:ch14-burst-chart}.

\noindent\textbf{Listing 14.3 --- \texttt{burst\_simulation.py}: seeded
replay harness feeding the comparison chart.}
\label{lst:ch14-sim}
\begin{minted}{python}
"""Simulation harness: replay one synthetic request timeline through every
reference limiter and print per-bucket admission counts.

The timeline is fully seeded and deterministic; re-running always prints the
same numbers. The printed coordinates are exactly the data series plotted in
Figure 14.4 (burst behavior comparison).

Scenario (Northwind Robotics telemetry ingest API, nominal plan:
10 requests/second per API key):
  * t in [0.90, 1.10]: a 25-request burst straddling the t=1.0 fixed-window
    boundary (a mobile fleet retrying in sync after a connectivity blip),
  * t in [1.50, 4.00]: a steady stream of 8 requests/second with seeded
    jitter -- well below the nominal rate, so a correct limiter should admit
    all of it after the burst penalty is paid.
"""

from __future__ import annotations

import random
from typing import Callable, Dict, List, Tuple

from rate_limiters import (
    FixedWindowLimiter,
    GCRALimiter,
    SlidingWindowCounterLimiter,
    SlidingWindowLogLimiter,
    TokenBucketLimiter,
    VirtualClock,
)

NOMINAL_RATE = 10  # requests per second per API key
WINDOW = 1.0
BUCKET_WIDTH = 0.25  # chart resolution


def build_timeline(seed: int = 14) -> List[float]:
    """Deterministic synthetic request timeline (arrival times, seconds)."""
    rng = random.Random(seed)
    arrivals: List[float] = []
    # Synchronized retry burst: 25 requests spread across the t=1.0 boundary.
    for i in range(25):
        arrivals.append(0.90 + i * 0.008 + rng.uniform(0.0, 0.002))
    # Steady-state traffic at 8 rps with jitter, from t=1.5 to t=4.0.
    t = 1.5
    while t <= 4.0:
        arrivals.append(t)
        t += 1.0 / 8.0 + rng.uniform(-0.03, 0.03)
    arrivals.sort()
    return arrivals


def simulate(arrivals: List[float],
             make_limiter: Callable[[VirtualClock], object]) -> Dict[float, int]:
    """Replay arrivals; return {bucket_start: admitted_count}."""
    clock = VirtualClock()
    limiter = make_limiter(clock)
    buckets: Dict[float, int] = {}
    for t in arrivals:
        decision = limiter.check(at=t)  # type: ignore[attr-defined]
        if decision.allowed:
            bucket = round(int(t / BUCKET_WIDTH) * BUCKET_WIDTH, 2)
            buckets[bucket] = buckets.get(bucket, 0) + 1
    return buckets


def main() -> None:
    arrivals = build_timeline()
    factories: List[Tuple[str, Callable[[VirtualClock], object]]] = [
        ("fixed window", lambda c: FixedWindowLimiter(NOMINAL_RATE, WINDOW, c)),
        ("sliding log", lambda c: SlidingWindowLogLimiter(NOMINAL_RATE, WINDOW, c)),
        ("sliding counter", lambda c: SlidingWindowCounterLimiter(NOMINAL_RATE, WINDOW, c)),
        ("token bucket", lambda c: TokenBucketLimiter(NOMINAL_RATE, NOMINAL_RATE, c)),
        ("GCRA", lambda c: GCRALimiter(NOMINAL_RATE, NOMINAL_RATE, c)),
    ]
    print(f"timeline: {len(arrivals)} arrivals, "
          f"t={arrivals[0]:.3f}..{arrivals[-1]:.3f}")
    all_starts = [round(i * BUCKET_WIDTH, 2) for i in range(int(4.0 / BUCKET_WIDTH) + 1)]
    for name, factory in factories:
        buckets = simulate(arrivals, factory)
        total = sum(buckets.values())
        coords = " ".join(f"({s:.2f}, {buckets.get(s, 0)})" for s in all_starts)
        print(f"\n{name}: admitted {total}/{len(arrivals)}")
        print(coords)


if __name__ == "__main__":
    main()
\end{minted}

\begin{figure}[ht]
  \centering
  \begin{tikzpicture}
    \begin{axis}[
        width=0.9\textwidth,
        height=7.6cm,
        xlabel={Time (s)},
        ylabel={Admitted requests per 250 ms bucket},
        xmin=0.5, xmax=4.1,
        ymin=0, ymax=21,
        legend pos=north east,
        legend style={font=\footnotesize},
        grid=major,
        grid style={gray!20},
      ]
      \addplot[only marks, mark=*, color=packtorange, mark size=1.6pt] coordinates {
        (0.75, 10) (1.00, 10) (1.25, 0) (1.50, 0) (1.75, 0) (2.00, 2)
        (2.25, 1) (2.50, 2) (2.75, 2) (3.00, 2) (3.25, 2) (3.50, 3) (3.75, 2)
      };
      \addlegendentry{Fixed window (36 admitted)}
      \addplot[only marks, mark=square*, color=darkblue, mark size=1.6pt] coordinates {
        (0.75, 10) (1.00, 0) (1.25, 0) (1.50, 0) (1.75, 1) (2.00, 2)
        (2.25, 1) (2.50, 2) (2.75, 2) (3.00, 2) (3.25, 2) (3.50, 3) (3.75, 2)
      };
      \addlegendentry{Sliding log (27 admitted)}
      \addplot[only marks, mark=triangle*, color=green!50!black, mark size=1.9pt] coordinates {
        (0.75, 10) (1.00, 1) (1.25, 0) (1.50, 3) (1.75, 2) (2.00, 2)
        (2.25, 1) (2.50, 2) (2.75, 2) (3.00, 2) (3.25, 2) (3.50, 3) (3.75, 2)
      };
      \addlegendentry{Sliding counter (32 admitted)}
      \addplot[only marks, mark=diamond*, color=purple, mark size=1.9pt] coordinates {
        (0.75, 10) (1.00, 1) (1.25, 0) (1.50, 3) (1.75, 2) (2.00, 2)
        (2.25, 1) (2.50, 2) (2.75, 2) (3.00, 2) (3.25, 2) (3.50, 3) (3.75, 2)
      };
      \addlegendentry{Token bucket = GCRA (32 admitted)}
    \end{axis}
  \end{tikzpicture}
  \caption{Burst behavior at the same nominal rate (10 requests/second per
  key), measured by Listing~\ref{lst:ch14-sim}. The fixed window admits 20
  requests across the $t{=}1.0$ boundary --- Theorem~\ref{thm:ch14-fixed-window-2x}
  made visible --- while the sliding log pays the burst penalty back in
  full before resuming. The token bucket and GCRA series coincide point for
  point, as Proposition~\ref{prop:ch14-gcra-equiv} requires; the sliding
  counter's weighted estimate coincidentally matches them on this timeline.}
  \label{fig:ch14-burst-chart}
\end{figure}
\subsection{Listing 14.4 --- Distributed Sliding Window: Redis + Lua}
\label{subsec:ch14-listing-redis}

The distributed limiter lifts the sliding log of
Section~\ref{subsec:ch14-sliding-log} into Redis sorted sets, with the
entire evict--count--admit--expire sequence inside one Lua script so it is
atomic by construction (Section~\ref{subsec:ch14-distributed}). The script
first, because it is the load-bearing artifact:

\noindent\textbf{Listing 14.4a --- \texttt{sliding\_window.lua}: atomic
check-and-admit (evict, count, conditionally insert, expire; returns
admission, remaining budget, and retry-after in milliseconds).}
\label{lst:ch14-redis-lua}
\begin{minted}{lua}
local key     = KEYS[1]
local now     = tonumber(ARGV[1])
local window  = tonumber(ARGV[2])
local limit   = tonumber(ARGV[3])
local member  = ARGV[4]

redis.call('ZREMRANGEBYSCORE', key, '-inf', now - window)
local count = redis.call('ZCARD', key)

if count < limit then
  redis.call('ZADD', key, now, member)
  redis.call('PEXPIRE', key, math.ceil(window * 1000))
  return {1, limit - count - 1, 0}
else
  -- ZRANGE + ZSCORE (not WITHSCORES): the WITHSCORES return shape inside
  -- Lua differs between Redis and emulators; two calls are portable.
  local oldest = redis.call('ZRANGE', key, 0, 0)
  local retry_after_ms = 0
  if #oldest >= 1 then
    local score = tonumber(redis.call('ZSCORE', key, oldest[1]))
    retry_after_ms = math.ceil((score + window - now) * 1000)
  end
  return {0, 0, retry_after_ms}
end
\end{minted}

Two script details earn their comments. First, Lua return values cross back
into RESP as \emph{integers}, so sub-second \texttt{Retry-After} precision
must travel as milliseconds and be divided down client-side. Second, the
member must be unique per request --- \texttt{ZADD} upserts by member, so
two requests sharing a member would silently undercount. The Python wrapper
guarantees uniqueness with a lock-guarded sequence counter.

\noindent\textbf{Listing 14.4b --- \texttt{redis\_limiter.py}: the wrapper,
plus an offline fakeredis factory.}
\label{lst:ch14-redis}
\begin{minted}{python}
"""Distributed sliding-window limiter: Redis sorted set + Lua script.

State lives in Redis so any number of API replicas enforce one shared limit
per key. The whole check-and-admit operation is a single Lua script, so it
executes atomically on the Redis server: no read-modify-write race can split
it, no matter how many replicas call concurrently.

Offline-first: tests run against ``fakeredis.FakeRedis``, which emulates
RESP and executes Lua scripts atomically inside the process, so the same
code paths are exercised without a server. For a real deployment swap in
``redis.Redis`` (redis>=5) -- the script and calling convention are
identical.
"""

from __future__ import annotations

import itertools
import threading
from typing import Optional, Protocol

from rate_limiters import Clock, Decision, VirtualClock


class RedisLike(Protocol):
    """The subset of the redis-py API this limiter needs."""

    def eval(self, script: str, numkeys: int, *keys_and_args: object) -> object: ...


# Sliding-window log as one atomic Lua script (also printed standalone as
# Listing 14.4a; it lives inline here so this file is complete and runnable).
SLIDING_WINDOW_LUA = """
local key     = KEYS[1]
local now     = tonumber(ARGV[1])
local window  = tonumber(ARGV[2])
local limit   = tonumber(ARGV[3])
local member  = ARGV[4]

redis.call('ZREMRANGEBYSCORE', key, '-inf', now - window)
local count = redis.call('ZCARD', key)

if count < limit then
  redis.call('ZADD', key, now, member)
  redis.call('PEXPIRE', key, math.ceil(window * 1000))
  return {1, limit - count - 1, 0}
else
  -- ZRANGE + ZSCORE (not WITHSCORES): the WITHSCORES return shape inside
  -- Lua differs between Redis and emulators; two calls are portable.
  local oldest = redis.call('ZRANGE', key, 0, 0)
  local retry_after_ms = 0
  if #oldest >= 1 then
    local score = tonumber(redis.call('ZSCORE', key, oldest[1]))
    retry_after_ms = math.ceil((score + window - now) * 1000)
  end
  return {0, 0, retry_after_ms}
end
"""


class RedisSlidingLogLimiter:
    """Shared sliding-window-log limiter backed by Redis sorted sets.

    Semantics match ``SlidingWindowLogLimiter`` exactly: at most ``limit``
    admissions per key inside any interval of ``window`` seconds. One Redis
    key per (namespace, limiter key) pair; keys self-expire after ``window``
    so idle consumers cost no memory.
    """

    def __init__(
        self,
        client: RedisLike,
        namespace: str,
        limit: int,
        window: float,
        clock: Clock,
    ) -> None:
        if limit <= 0 or window <= 0:
            raise ValueError("limit and window must be positive")
        if not namespace:
            raise ValueError("namespace must be non-empty")
        self.client = client
        self.namespace = namespace
        self.limit = limit
        self.window = float(window)
        self.clock = clock
        self._seq = itertools.count()
        self._seq_lock = threading.Lock()

    def _member(self, t: float) -> str:
        # ZADD upserts by member, so every admitted request needs a unique
        # member id; identical scores are fine.
        with self._seq_lock:
            return f"{t:.6f}:{next(self._seq)}"

    def check(self, key: str, at: Optional[float] = None) -> Decision:
        if not key:
            raise ValueError("limiter key must be non-empty")
        t = self.clock.now() if at is None else at
        redis_key = f"{self.namespace}:{key}"
        raw = self.client.eval(
            SLIDING_WINDOW_LUA,
            1,
            redis_key,
            repr(t),
            repr(self.window),
            str(self.limit),
            self._member(t),
        )
        allowed, remaining, retry_after_ms = (int(v) for v in raw)
        return Decision(bool(allowed), remaining, retry_after_ms / 1000.0)


def build_fake_limiter(
    limit: int = 5, window: float = 10.0, clock: Optional[Clock] = None
) -> RedisSlidingLogLimiter:
    """Offline factory: a limiter against an in-process fakeredis server."""
    import fakeredis

    return RedisSlidingLogLimiter(
        client=fakeredis.FakeRedis(),
        namespace="northwind:rl",
        limit=limit,
        window=window,
        clock=clock or VirtualClock(),
    )
\end{minted}

\noindent\textbf{Listing 14.5 --- \texttt{test\_redis\_limiter.py}:
semantics, isolation, expiry hygiene, and a 16-thread concurrency storm.}
\label{lst:ch14-redis-tests}
\begin{minted}{python}
"""Tests for the Redis+Lua distributed limiter (offline, fakeredis).

Atomicity note: real Redis executes each Lua script atomically on its single
command thread, so no interleaving of ZREMRANGEBYSCORE/ZCARD/ZADD is
possible across clients. fakeredis emulates the same contract in-process
(scripts run to completion under its server lock), so the concurrency test
below exercises the true atomicity semantics without a server. What fakeredis
does NOT emulate is wall-clock key expiry timing -- which is why the limiter
never relies on EXPIRE for correctness, only for memory hygiene.
"""

from __future__ import annotations

import threading

import pytest

from rate_limiters import VirtualClock
from redis_limiter import RedisSlidingLogLimiter, build_fake_limiter


class TestSemantics:
    def test_admits_limit_then_rejects(self) -> None:
        clock = VirtualClock()
        limiter = build_fake_limiter(limit=3, window=10.0, clock=clock)
        results = [limiter.check("robot-7", at=1.0).allowed for _ in range(5)]
        assert results == [True, True, True, False, False]

    def test_window_slides_exactly(self) -> None:
        clock = VirtualClock()
        limiter = build_fake_limiter(limit=2, window=10.0, clock=clock)
        limiter.check("robot-7", at=5.0)
        limiter.check("robot-7", at=7.0)
        assert not limiter.check("robot-7", at=14.999).allowed
        assert limiter.check("robot-7", at=15.0).allowed  # 5.0 has aged out

    def test_keys_are_isolated(self) -> None:
        clock = VirtualClock()
        limiter = build_fake_limiter(limit=1, window=10.0, clock=clock)
        assert limiter.check("robot-7", at=0.0).allowed
        assert limiter.check("robot-8", at=0.0).allowed
        assert not limiter.check("robot-7", at=0.0).allowed

    def test_retry_after_matches_oldest_expiry(self) -> None:
        clock = VirtualClock()
        limiter = build_fake_limiter(limit=2, window=10.0, clock=clock)
        limiter.check("robot-7", at=1.25)
        limiter.check("robot-7", at=3.00)
        decision = limiter.check("robot-7", at=4.0)
        assert not decision.allowed
        # Oldest (1.25) leaves the window at 11.25; ceil to milliseconds.
        assert decision.retry_after == pytest.approx(7.25, abs=0.001)

    def test_idle_keys_expire_from_redis(self) -> None:
        import fakeredis

        clock = VirtualClock()
        client = fakeredis.FakeRedis()
        limiter = RedisSlidingLogLimiter(client, "northwind:rl", 5, 10.0, clock)
        limiter.check("robot-7", at=0.0)
        ttl_ms = client.pttl("northwind:rl:robot-7")
        assert 9_000 < ttl_ms <= 10_000


class TestConcurrency:
    def test_no_over_admission_under_thread_storm(self) -> None:
        """16 threads x 25 requests against limit 100: exactly 100 admissions,
        never 101. A non-atomic check-then-set would intermittently fail."""
        clock = VirtualClock()
        limiter = build_fake_limiter(limit=100, window=60.0, clock=clock)
        admitted = 0
        admitted_lock = threading.Lock()

        def worker() -> None:
            nonlocal admitted
            local = 0
            for _ in range(25):
                if limiter.check("fleet-shared-key", at=0.0).allowed:
                    local += 1
            with admitted_lock:
                admitted += local

        threads = [threading.Thread(target=worker) for _ in range(16)]
        for thread in threads:
            thread.start()
        for thread in threads:
            thread.join()
        assert admitted == 100
\end{minted}

\subsection{Listing 14.6 --- FastAPI Middleware: Headers Are the Product}
\label{subsec:ch14-listing-middleware}

The middleware implements the full client-facing contract of
Section~\ref{subsec:ch14-headers}: per-API-key buckets, 429 with ceiled
delay-seconds \texttt{Retry-After}, the four IETF \texttt{RateLimit-*}
fields on successes \emph{and} rejections, 401 for missing keys (distinct
from 429 --- an unauthenticated caller has no bucket to exceed), and a
fail-open/fail-closed policy switch for limiter outages. It is written as
raw ASGI middleware rather than \texttt{BaseHTTPMiddleware} so that rejected
requests are answered without ever instantiating the downstream
application's request handling.

\noindent\textbf{Listing 14.6 --- \texttt{middleware\_app.py}: rate-limit
middleware with accurate headers and a tested failure policy.}
\label{lst:ch14-middleware}
\begin{minted}{python}
"""FastAPI middleware wiring a rate limiter to client-facing HTTP semantics.

Semantics implemented (Section 14.3, "Client-Facing Semantics"):
  * per-API-key buckets, keyed on the ``X-API-Key`` header,
  * 429 Too Many Requests with ``Retry-After`` (delay-seconds form) on
    rejection,
  * IETF RateLimit header fields on every limited response:
    ``RateLimit-Limit``, ``RateLimit-Remaining``, ``RateLimit-Reset``,
    ``RateLimit-Policy`` (e.g. ``5;w=60``),
  * a fail-open / fail-closed policy switch for limiter outages.

The limiter is injected behind the ``KeyedLimiter`` protocol, so the same
middleware runs against the in-process fixed-window limiter here, or the
Redis limiter of Listing 14.4 in a multi-replica deployment. The clock is
injected, so tests are deterministic.
"""

from __future__ import annotations

import math
from typing import Dict, Literal, Optional, Protocol

from fastapi import FastAPI, Request
from fastapi.responses import JSONResponse
from fastapi.testclient import TestClient

from rate_limiters import Clock, Decision, VirtualClock

FailurePolicy = Literal["fail_open", "fail_closed"]

VALID_API_KEYS = {"nr-demo-key-001", "nr-demo-key-002"}  # synthetic only


class KeyedLimiter(Protocol):
    """Anything that can admit or reject for a key and report reset time."""

    def check(self, key: str) -> Decision: ...

    def reset_after(self, key: str) -> float:
        """Seconds until the key's window fully resets (0 if idle)."""
        ...


class InProcessFixedWindowLimiter:
    """Per-key fixed-window counters with an injected clock.

    Correct O(1)-memory semantics for a single replica; subject to the
    2x boundary burst of Theorem 14.1, which is acceptable when the window
    is short and the header contract (crisp RateLimit-Reset) matters more
    than boundary strictness. Swap in RedisSlidingLogLimiter for
    multi-replica deployments.
    """

    def __init__(self, limit: int, window: float, clock: Clock) -> None:
        if limit <= 0 or window <= 0:
            raise ValueError("limit and window must be positive")
        self.limit = limit
        self.window = window
        self.clock = clock
        self._buckets: Dict[str, tuple[int, int]] = {}  # key -> (index, count)

    def check(self, key: str) -> Decision:
        t = self.clock.now()
        index = math.floor(t / self.window)
        bucket_index, count = self._buckets.get(key, (index, 0))
        if bucket_index != index:
            bucket_index, count = index, 0
        window_end = (index + 1) * self.window
        if count < self.limit:
            self._buckets[key] = (bucket_index, count + 1)
            return Decision(True, self.limit - count - 1, 0.0)
        self._buckets[key] = (bucket_index, count)
        return Decision(False, 0, window_end - t)

    def reset_after(self, key: str) -> float:
        t = self.clock.now()
        index = math.floor(t / self.window)
        bucket = self._buckets.get(key)
        if bucket is None or bucket[0] != index or bucket[1] == 0:
            return 0.0
        return (index + 1) * self.window - t


class RateLimitMiddleware:
    """ASGI middleware enforcing a per-API-key limit with full headers."""

    def __init__(
        self,
        app: FastAPI,
        limiter: KeyedLimiter,
        limit: int,
        window: float,
        on_limiter_error: FailurePolicy = "fail_open",
    ) -> None:
        if on_limiter_error not in ("fail_open", "fail_closed"):
            raise ValueError("on_limiter_error must be 'fail_open' or 'fail_closed'")
        self.app = app
        self.limiter = limiter
        self.limit = limit
        self.window = window
        self.on_limiter_error = on_limiter_error

    async def __call__(self, scope, receive, send) -> None:
        if scope["type"] != "http":
            await self.app(scope, receive, send)
            return
        request = Request(scope)
        api_key = request.headers.get("x-api-key")
        if api_key not in VALID_API_KEYS:
            response = JSONResponse(
                status_code=401,
                content={"detail": "missing or invalid API key"},
            )
            await response(scope, receive, send)
            return

        try:
            decision = self.limiter.check(api_key)
        except Exception:
            if self.on_limiter_error == "fail_open":
                # Limiter down: serve anyway, and say nothing we cannot
                # substantiate -- no RateLimit headers we cannot compute.
                await self.app(scope, receive, send)
                return
            response = JSONResponse(
                status_code=503,
                content={"detail": "rate limiter unavailable; failing closed"},
                headers={"Retry-After": "5"},
            )
            await response(scope, receive, send)
            return

        headers = self._ratelimit_headers(api_key, decision)
        if not decision.allowed:
            headers["Retry-After"] = str(max(1, math.ceil(decision.retry_after)))
            response = JSONResponse(
                status_code=429,
                content={"detail": "rate limit exceeded"},
                headers=headers,
            )
            await response(scope, receive, send)
            return

        async def send_with_headers(message) -> None:
            if message["type"] == "http.response.start":
                extra = [(k.lower().encode(), v.encode()) for k, v in headers.items()]
                message["headers"] = list(message.get("headers", [])) + extra
            await send(message)

        await self.app(scope, receive, send_with_headers)

    def _ratelimit_headers(self, key: str, decision: Decision) -> Dict[str, str]:
        reset = decision.retry_after if not decision.allowed else self.limiter.reset_after(key)
        return {
            "RateLimit-Limit": str(self.limit),
            "RateLimit-Remaining": str(max(0, decision.remaining)),
            "RateLimit-Reset": str(max(0, math.ceil(reset))),
            "RateLimit-Policy": f"{self.limit};w={int(self.window)}",
        }


def create_app(
    limiter: Optional[KeyedLimiter] = None,
    clock: Optional[Clock] = None,
    on_limiter_error: FailurePolicy = "fail_open",
    limit: int = 5,
    window: float = 60.0,
) -> FastAPI:
    """Northwind Robotics telemetry API behind the rate-limit middleware."""
    clock = clock or VirtualClock()
    limiter = limiter or InProcessFixedWindowLimiter(limit, window, clock)
    app = FastAPI(title="Northwind Robotics Telemetry API")
    app.add_middleware(
        RateLimitMiddleware,
        limiter=limiter,
        limit=limit,
        window=window,
        on_limiter_error=on_limiter_error,
    )

    @app.get("/telemetry")
    def telemetry() -> dict[str, str]:
        return {"status": "ok", "fleet": "northwind-robotics"}

    return app


if __name__ == "__main__":
    demo_clock = VirtualClock()
    client = TestClient(create_app(clock=demo_clock))
    for i in range(7):
        resp = client.get("/telemetry", headers={"X-API-Key": "nr-demo-key-001"})
        interesting = {k: v for k, v in resp.headers.items()
                       if k.startswith(("ratelimit", "retry"))}
        print(f"request {i + 1}: {resp.status_code} {interesting}")
\end{minted}

\noindent\textbf{Listing 14.7 --- \texttt{test\_middleware\_app.py}: header
accuracy across window boundaries and both failure policies.}
\label{lst:ch14-middleware-tests}
\begin{minted}{python}
"""TestClient tests: header accuracy across window boundaries, 429 +
Retry-After, and the fail-open / fail-closed policy switch. All offline,
all deterministic (injected VirtualClock).
"""

from __future__ import annotations

import pytest
from fastapi.testclient import TestClient

from middleware_app import create_app
from rate_limiters import Decision, VirtualClock

KEY = {"X-API-Key": "nr-demo-key-001"}
OTHER_KEY = {"X-API-Key": "nr-demo-key-002"}


@pytest.fixture()
def rig() -> tuple[TestClient, VirtualClock]:
    clock = VirtualClock(start=0.0)
    app = create_app(clock=clock, limit=5, window=60.0)
    return TestClient(app), clock


class TestHeaderAccuracy:
    def test_limit_remaining_and_policy_on_success(self, rig) -> None:
        client, _ = rig
        for expected_remaining in (4, 3, 2, 1, 0):
            resp = client.get("/telemetry", headers=KEY)
            assert resp.status_code == 200
            assert resp.headers["RateLimit-Limit"] == "5"
            assert resp.headers["RateLimit-Remaining"] == str(expected_remaining)
            assert resp.headers["RateLimit-Policy"] == "5;w=60"

    def test_reset_counts_down_to_window_end(self, rig) -> None:
        client, clock = rig
        resp = client.get("/telemetry", headers=KEY)
        assert resp.headers["RateLimit-Reset"] == "60"
        clock.advance(25.0)
        resp = client.get("/telemetry", headers=KEY)
        assert resp.headers["RateLimit-Reset"] == "35"  # 60 - 25

    def test_429_carries_retry_after_and_zero_remaining(self, rig) -> None:
        client, clock = rig
        for _ in range(5):
            assert client.get("/telemetry", headers=KEY).status_code == 200
        clock.advance(10.5)  # 49.5s until window end
        resp = client.get("/telemetry", headers=KEY)
        assert resp.status_code == 429
        assert resp.headers["Retry-After"] == "50"  # ceil(49.5)
        assert resp.headers["RateLimit-Remaining"] == "0"
        assert resp.headers["RateLimit-Reset"] == "50"

    def test_headers_reset_across_window_boundary(self, rig) -> None:
        client, clock = rig
        for _ in range(5):
            client.get("/telemetry", headers=KEY)
        assert client.get("/telemetry", headers=KEY).status_code == 429
        clock.advance(60.0)  # cross the boundary at t=60
        resp = client.get("/telemetry", headers=KEY)
        assert resp.status_code == 200
        assert resp.headers["RateLimit-Remaining"] == "4"
        assert resp.headers["RateLimit-Reset"] == "60"

    def test_buckets_are_per_api_key(self, rig) -> None:
        client, _ = rig
        for _ in range(5):
            client.get("/telemetry", headers=KEY)
        assert client.get("/telemetry", headers=KEY).status_code == 429
        resp = client.get("/telemetry", headers=OTHER_KEY)
        assert resp.status_code == 200
        assert resp.headers["RateLimit-Remaining"] == "4"

    def test_missing_key_is_401_not_429(self, rig) -> None:
        client, _ = rig
        resp = client.get("/telemetry")
        assert resp.status_code == 401
        assert "RateLimit-Limit" not in resp.headers


class BrokenLimiter:
    """A limiter that is down: every check raises, like a lost Redis."""

    def check(self, key: str) -> Decision:
        raise ConnectionError("redis: connection refused")

    def reset_after(self, key: str) -> float:
        raise ConnectionError("redis: connection refused")


class TestFailurePolicy:
    def test_fail_open_serves_without_unverifiable_headers(self) -> None:
        app = create_app(limiter=BrokenLimiter(), on_limiter_error="fail_open")
        client = TestClient(app)
        resp = client.get("/telemetry", headers=KEY)
        assert resp.status_code == 200
        assert "RateLimit-Remaining" not in resp.headers

    def test_fail_closed_returns_503_with_retry_after(self) -> None:
        app = create_app(limiter=BrokenLimiter(), on_limiter_error="fail_closed")
        client = TestClient(app)
        resp = client.get("/telemetry", headers=KEY)
        assert resp.status_code == 503
        assert resp.headers["Retry-After"] == "5"
\end{minted}

\subsection{Listing 14.8 --- Cost-Based Limiting: Budgeting Work, Not Requests}
\label{subsec:ch14-listing-cost}

The cost-based limiter of Section~\ref{subsec:ch14-multidim} is a token
bucket whose \texttt{cost} parameter is looked up per endpoint: the budget
is denominated in normalized work units (server CPU-milliseconds, from the
observability pipeline), so an export drains it ten times faster than a
median read, and \texttt{/health} is free because throttling a liveness
probe is how you get load balancers to shoot healthy servers.

\noindent\textbf{Listing 14.8 --- \texttt{cost\_limiter.py}: per-endpoint
weighted token buckets.}
\label{lst:ch14-cost}
\begin{minted}{python}
"""Cost-based (weighted) rate limiting: one token bucket per API key, with
each endpoint charging a weight proportional to its true server-side cost.

A flat "requests per minute" limit treats a ``GET /health`` and a
``POST /reports/export`` as equal, but the export scans and renders a
million-row dataset while the health check is a no-op. Weighted buckets make
the budget denominated in *work*, not in *requests*: expensive operations
drain the shared budget faster, free endpoints (health probes) are never
throttled, and cheap endpoints survive exactly as long as any budget
remains.

Northwind Robotics cost table: weights are server CPU-millisecond estimates
from the Chapter 18 observability pipeline, normalized so the median read
endpoint costs 1.0.
"""

from __future__ import annotations

from typing import Dict, Optional

from rate_limiters import Clock, Decision, TokenBucketLimiter, VirtualClock

# path -> token cost per call (normalized work units)
ENDPOINT_COSTS: Dict[str, float] = {
    "/health": 0.0,          # liveness probe: never throttle
    "/telemetry": 1.0,       # median read
    "/fleet/status": 1.0,    # median read
    "/telemetry/batch": 4.0, # bulk ingest: 4x a single write
    "/reports/export": 10.0, # full-fleet CSV render: the expensive one
}
DEFAULT_COST = 1.0


class CostBasedLimiter:
    """Per-key token buckets charged by endpoint weight.

    Budget semantics: ``rate`` work-units per second sustained, ``burst``
    work-units of instantaneous headroom. With rate=10 and burst=20, a key
    may run 2 exports instantly (20 / 10) or 20 median reads, or any mix
    whose total weight fits the bucket.
    """

    def __init__(self, rate: float, burst: float, clock: Clock,
                 costs: Optional[Dict[str, float]] = None) -> None:
        if rate <= 0 or burst <= 0:
            raise ValueError("rate and burst must be positive")
        self.rate = rate
        self.burst = burst
        self.clock = clock
        self.costs = dict(ENDPOINT_COSTS if costs is None else costs)
        self._buckets: Dict[str, TokenBucketLimiter] = {}

    def cost_of(self, path: str) -> float:
        return self.costs.get(path, DEFAULT_COST)

    def check(self, key: str, path: str, at: Optional[float] = None) -> Decision:
        if not key:
            raise ValueError("limiter key must be non-empty")
        cost = self.cost_of(path)
        if cost == 0.0:
            return Decision(True, self.limit_placeholder(), 0.0)
        bucket = self._buckets.get(key)
        if bucket is None:
            bucket = TokenBucketLimiter(self.rate, self.burst, self.clock)
            self._buckets[key] = bucket
        return bucket.check(at=at, cost=cost)

    def limit_placeholder(self) -> int:
        # Free endpoints have no meaningful "remaining"; report the burst
        # budget so header math stays well-formed.
        return int(self.burst)


if __name__ == "__main__":
    clock = VirtualClock()
    limiter = CostBasedLimiter(rate=10.0, burst=20.0, clock=clock)
    key = "nr-demo-key-001"
    for path in ("/reports/export", "/reports/export", "/reports/export",
                 "/telemetry", "/telemetry", "/health"):
        d = limiter.check(key, path)
        verdict = "admitted" if d.allowed else f"rejected (retry in {d.retry_after:.1f}s)"
        print(f"{path:<18} cost={limiter.cost_of(path):>4} -> {verdict}")
\end{minted}

\noindent\textbf{Listing 14.9 --- \texttt{test\_cost\_limiter.py}: weighted
drainage, shared budgets, and edge cases.}
\label{lst:ch14-cost-tests}
\begin{minted}{python}
"""Tests for the cost-based limiter: weighted drainage, free endpoints,
default costs, and graceful degradation of cheap endpoints under budget
pressure. Deterministic via VirtualClock.
"""

from __future__ import annotations

import pytest

from cost_limiter import CostBasedLimiter
from rate_limiters import VirtualClock

KEY = "nr-demo-key-001"


@pytest.fixture()
def limiter() -> CostBasedLimiter:
    return CostBasedLimiter(rate=10.0, burst=20.0, clock=VirtualClock())


class TestWeightedDrainage:
    def test_export_costs_ten_reads(self, limiter: CostBasedLimiter) -> None:
        # Two exports (2 x 10 units) exactly exhaust a 20-unit bucket.
        assert limiter.check(KEY, "/reports/export", at=0.0).allowed
        assert limiter.check(KEY, "/reports/export", at=0.0).allowed
        assert not limiter.check(KEY, "/reports/export", at=0.0).allowed

    def test_mixed_workload_shares_one_budget(self, limiter: CostBasedLimiter) -> None:
        assert limiter.check(KEY, "/reports/export", at=0.0).allowed  # 10
        for _ in range(10):
            assert limiter.check(KEY, "/telemetry", at=0.0).allowed   # 10 x 1
        assert not limiter.check(KEY, "/telemetry", at=0.0).allowed   # 21st unit

    def test_refill_is_in_work_units(self, limiter: CostBasedLimiter) -> None:
        limiter.check(KEY, "/reports/export", at=0.0)
        limiter.check(KEY, "/reports/export", at=0.0)
        # 1.0s of refill = 10 work units = exactly one more export.
        assert limiter.check(KEY, "/reports/export", at=1.0).allowed
        assert not limiter.check(KEY, "/reports/export", at=1.0).allowed

    def test_retry_after_scales_with_cost(self, limiter: CostBasedLimiter) -> None:
        limiter.check(KEY, "/reports/export", at=0.0)
        limiter.check(KEY, "/reports/export", at=0.0)
        cheap = limiter.check(KEY, "/telemetry", at=0.0)
        dear = limiter.check(KEY, "/reports/export", at=0.0)
        assert cheap.retry_after == pytest.approx(0.1)   # 1 unit / 10 rps
        assert dear.retry_after == pytest.approx(1.0)    # 10 units / 10 rps


class TestEdgeCases:
    def test_health_endpoint_is_never_throttled(self, limiter: CostBasedLimiter) -> None:
        for _ in range(2):
            limiter.check(KEY, "/reports/export", at=0.0)
        for _ in range(100):
            assert limiter.check(KEY, "/health", at=0.0).allowed

    def test_unknown_path_defaults_to_unit_cost(self, limiter: CostBasedLimiter) -> None:
        assert limiter.cost_of("/no/such/route") == 1.0
        for i in range(20):
            assert limiter.check(KEY, "/no/such/route", at=0.0).allowed, i
        assert not limiter.check(KEY, "/no/such/route", at=0.0).allowed

    def test_budgets_are_per_key(self, limiter: CostBasedLimiter) -> None:
        limiter.check(KEY, "/reports/export", at=0.0)
        limiter.check(KEY, "/reports/export", at=0.0)
        assert not limiter.check(KEY, "/reports/export", at=0.0).allowed
        assert limiter.check("nr-demo-key-002", "/reports/export", at=0.0).allowed
\end{minted}

The complete suite --- 35 tests across the four test listings --- runs in
about four seconds with no network, no containers, and no wall-clock
dependence:

\begin{minted}{text}
pytest -q
# 35 passed
python burst_simulation.py   # prints the Figure 14.4 coordinates
python middleware_app.py     # 7-request demo: headers decay, then 429
python cost_limiter.py       # weighted drainage demo
\end{minted}
%% ----------------------------------------------------------------------------
\section{Common Pitfalls \& Anti-Patterns}
\label{sec:ch14-pitfalls}

\begin{enumerate}
  \item \textbf{Fixed window at scale, believed to be strict.} The 2x
        boundary burst (Theorem~\ref{thm:ch14-fixed-window-2x}) is
        disqualifying when the limit protects something with a hard
        per-second ceiling (a partner API, a connection pool). If you keep
        the fixed window for its crisp reset headers, halve the nominal
        limit or shorten the window, and \emph{document the envelope you
        actually enforce}.
  \item \textbf{Per-IP as the only dimension.} Behind corporate NATs and
        carrier-grade NAT, one source IP can be a ten-thousand-employee
        enterprise: a tight per-IP limit punishes your best customers
        collectively for their network topology, while a botnet of ten
        thousand IPs gets ten thousand times the limit. Per-IP is a coarse
        outer net for pre-authentication traffic only; identity-based
        limits (API key, account, session) must carry the real weight ---
        Section~\ref{sec:ch14-case} is the cautionary tale in full.
  \item \textbf{429 without \texttt{Retry-After}.} A rejection that does not
        say when to come back forces clients to guess, and the industry's
        guess is immediate retry --- your limiter just built itself a retry
        storm. Always emit \texttt{Retry-After}; prefer delay-seconds;
        always ceil~\cite{rfc9110}.
  \item \textbf{Non-atomic check-then-set across replicas.}
        \texttt{GET}-compare-\texttt{INCR} over-admits under concurrency;
        \texttt{INCR}-then-reject lets rejected traffic poison the counter.
        One Lua script, one authoritative cell, or accept a documented
        approximation --- there is no fourth option
        (Section~\ref{subsec:ch14-distributed}).
  \item \textbf{Silent throttling.} Queueing or delaying requests without
        any observable signal (no headers, no logs the client can see, no
        documented policy) makes your API look randomly slow and puts the
        debugging cost on the wrong team. If you shape, say so.
  \item \textbf{One global limit for all endpoints.} A single
        requests-per-minute number treats \texttt{GET /health} and
        \texttt{POST /reports/export} as identical. Attackers and sloppy
        integrators will find your most expensive route and spend the whole
        budget there. Weight by cost (Listing~\ref{lst:ch14-cost}) and cap
        concurrency on the heavy routes separately.
  \item \textbf{Fail-closed limiter with no policy when Redis dies.}
        Deciding the failure posture \emph{during} the outage means the
        default decides for you --- usually fail-closed-by-exception, i.e.,
        the limiter takes the whole API down with its state store. Pick per
        endpoint class from Table~\ref{tab:ch14-failmatrix}, exempt health
        checks, and test both paths (Listing~\ref{lst:ch14-middleware-tests}
        does).
  \item \textbf{Wall-clock-dependent tests.} \texttt{time.sleep(1.1)} in a
        test suite is a flake factory: CI runners stall, laptops sleep, and
        the failure mode is intermittent --- the worst kind. Inject a clock;
        advance it explicitly. Every listing in this chapter demonstrates
        the pattern.
\end{enumerate}

%% ----------------------------------------------------------------------------
\section{Hands-On Production Lab \& Real-World Case Study}
\label{sec:ch14-lab}

\subsection*{Lab: From Simulation to Concurrency Storm}

\textbf{Goal.} Verify, on your own machine, every load-bearing claim of this
chapter: the algorithms' burst envelopes, header accuracy under a live
middleware, and fairness under a thread storm.

\textbf{Setup.} Save Listings 14.1--14.9 into one directory, create the
virtual environment from Section~\ref{sec:ch14-objectives}, and install
\texttt{fakeredis}, \texttt{fastapi}, \texttt{httpx}, \texttt{pytest}.

\begin{enumerate}
  \item \textbf{Reproduce the proofs.} \texttt{pytest test\_rate\_limiters.py -q}
        --- confirm 14 passes, then open
        \texttt{test\_double\_burst\_at\_boundary} and change the window to
        \texttt{2.0}: the proof still holds (20 admissions, now inside a
        2~ms interval of a ``5/s'' limit). Explain why $2N$ did not change.
  \item \textbf{Run the simulation.} \texttt{python burst\_simulation.py} ---
        compare each algorithm's per-bucket admissions to
        Figure~\ref{fig:ch14-burst-chart}. Then change the burst to land
        \emph{entirely inside} one fixed window (\texttt{0.50 + i * 0.008})
        and observe the fixed window's boundary advantage disappear.
  \item \textbf{Exercise the middleware.} \texttt{python middleware\_app.py}
        shows seven requests: \texttt{RateLimit-Remaining} decaying
        $4 \to 0$, then 429 with \texttt{Retry-After: 60}. Now load it: run
        \texttt{pytest test\_middleware\_app.py -q}, then write a small
        loop that fires two interleaved API keys through the
        \texttt{TestClient} and assert neither key's rejections delay the
        other's admissions --- fairness, verified.
  \item \textbf{Storm the distributed limiter.}
        \texttt{pytest test\_redis\_limiter.py -q} runs 16 threads
        $\times$ 25 requests against a limit of 100 and asserts exactly 100
        admissions. Break it deliberately: replace the Lua script with
        separate \texttt{zcard}/\texttt{zadd} calls and re-run until you
        catch the over-admission --- the race of
        Section~\ref{subsec:ch14-distributed}, reproduced on demand.
  \item \textbf{Flip the policy switch.} Run both \texttt{BrokenLimiter}
        tests, then decide for each Northwind Robotics endpoint
        (\texttt{/telemetry}, \texttt{/reports/export}, \texttt{/health},
        a hypothetical \texttt{/auth/login}) which failure policy belongs on
        it, and defend the assignment using
        Table~\ref{tab:ch14-failmatrix}.
\end{enumerate}

\subsection{Case Study: The Login Endpoint the Per-IP Limit Could Not Save}
\label{sec:ch14-case}

Northwind Robotics' fleet console had a textbook-looking defense:
\texttt{POST /auth/login} limited to 20 requests per minute per source IP,
fixed window, enforced at the edge. On a Tuesday, a credential-stuffing
campaign --- roughly 400{,}000 username/password pairs harvested from an
unrelated breach --- walked straight through it, and the postmortem found
\emph{three} compounding design errors rather than one.

\textbf{Error one: the limit dimension was the attacker's friend.} The
campaign rotated through a residential proxy pool of ~30{,}000 exit IPs, so
each IP made at most 13 attempts --- under the limit everywhere, forever.
The per-IP ceiling divided the attacker's budget by the size of their
botnet, which is to say it barely divided it at all.

\textbf{Error two: the same dimension was the customers' enemy.} Weeks
earlier, a 6{,}000-employee logistics customer behind a corporate NAT had
been \emph{collectively} locked out every morning at 09:00 when their
workforce logged in through two egress IPs. Support's fix --- raise the
per-IP limit to 2{,}000/min --- made error one cheaper to exploit. Per-IP
limits had now failed in both directions on the same endpoint.

\textbf{Error three: nothing counted the thing being attacked.} The asset
under attack was \emph{accounts}, but no counter was keyed by account.
There was no per-username attempt limit, no global login-endpoint budget,
and no velocity alert; the campaign's 0.7\% success rate against 400k pairs
was discovered by a customer reporting strange session emails, not by
telemetry.

The remediation layered the dimensions of
Section~\ref{subsec:ch14-multidim}: a per-account sliding log (5 attempts
per 10 minutes --- exactness chosen deliberately, $N$ tiny so $O(N)$ memory
is free) via the Redis cell of Listing~\ref{lst:ch14-redis}; a
\emph{global} per-endpoint token bucket on \texttt{/auth/login} as a
circuit-level budget; the per-IP limit retained but demoted to a loose
outer net (and fail-open, because a limiter outage must never lock out all
logins --- fail-closed is correct for the \emph{account} limiter instead);
exponential backoff plus progressive delay on repeated per-account
failures; and a velocity alert on global attempt rate. The follow-up
campaign a quarter later exhausted its per-account budgets inside ninety
seconds and the global bucket flagged the anomaly before support heard a
word. The lesson generalizes: \emph{choose the limiter dimension to match
the asset under attack, and assume the attacker chooses their own
distribution.}

%% ----------------------------------------------------------------------------
\section{Self-Assessment Exercises \& Deep-Dive Questions}
\label{sec:ch14-exercises}

\begin{enumerate}[label=\textbf{E14.\arabic*.}, leftmargin=2.2cm]
  \item \textbf{Boundary audit.} Modify Listing~\ref{lst:ch14-tests} to
        prove that the sliding counter, unlike the fixed window, admits
        \emph{fewer} than $2N$ requests when a burst straddles a boundary
        --- then exhibit a timeline where it still over-admits relative to
        the exact log. Which proposition bounds the difference?
  \item \textbf{Equivalence hunting.} Extend
        \texttt{test\_equivalent\_to\_token\_bucket} with 5{,}000 seeded
        events including multi-second idle gaps. Find (by search) a
        parameter pair $(r, B)$ where GCRA and the token bucket
        \emph{diverge}, or strengthen the proof sketch of
        Proposition~\ref{prop:ch14-gcra-equiv} to explain why none exists.
  \item \textbf{Reset semantics.} The middleware computes
        \texttt{RateLimit-Reset} from the fixed window boundary. Specify
        (and implement) what it should report for the \emph{sliding log}
        limiter, where ``the window'' never resets all at once. Which
        choice keeps the header monotonic and honest?
  \item \textbf{Clock discipline.} A teammate replaces the injected
        \texttt{Clock} with \texttt{time.monotonic} ``to simplify the
        API.'' List three concrete test or production failures this
        reintroduces, referencing specific listings.
  \item \textbf{Race reconstruction.} Instrument the broken
        non-atomic limiter from lab step~4 to \emph{count} over-admissions
        across 50 seeded storms, and plot (mentally or with pgfplots) how
        overshoot scales with thread count. Why does the Lua version's
        atomicity eliminate the entire class rather than reduce it?
  \item \textbf{Quota design.} Northwind Robotics adds a monthly export
        quota (100 exports) on top of the per-second cost budget. Design
        the storage key layout, the Lua-side interaction (one script or
        two?), and the \texttt{GET /v1/usage} response body. What happens
        to a client at 99/100 whose per-second bucket is also empty?
  \item \textbf{AIMD tuning.} Given additive increase $a = 5$ rps and
        decrease $\beta = 0.5$ evaluated once per second, with a service
        that overloads above 300 rps, simulate (on paper or in code) the
        limit trajectory from $L_0 = 100$ rps. How many seconds to first
        contact with capacity, and what is the steady-state oscillation
        amplitude? Why is $\beta < 1$ the safety-critical parameter rather
        than $a$?
  \item \textbf{Header archaeology.} Pick a public API you use. Capture a
        429 (or its headers near the limit) and document which of
        \texttt{Retry-After}, \texttt{RateLimit-Limit/-Remaining/-Reset},
        and \texttt{RateLimit-Policy} it emits, whether \texttt{Reset} is
        delay-seconds or epoch-seconds, and one concrete way its contract
        is better or worse than Listing~\ref{lst:ch14-middleware}'s.
\end{enumerate}
%% ----------------------------------------------------------------------------
\section{Interview Q\&A Vault}
\label{sec:ch14-interview}

\subsection*{Junior (Q1--Q10)}
\begin{enumerate}[label=\textbf{Q\arabic*.}, leftmargin=1.6cm]
  \item \textbf{What is rate limiting, in one paragraph?}
        Per-identity admission control over a short window: a rule like
        ``100 requests per minute per API key,'' enforced statefully, whose
        purpose is to protect shared capacity, keep tenants fair, control
        cost, and raise the price of abuse. Rejection (HTTP 429) is a
        normal, cheap, expected outcome --- not an error in the alarming
        sense.
  \item \textbf{Rate limiting versus throttling versus quota --- what's the
        difference?}
        Rate limiting \emph{rejects} excess starts in a short window;
        throttling \emph{shapes} flow by delaying or dequeuing admitted
        work; a quota is a long-window (hour/day/month) budget that rations
        total consumption, usually a business/plan construct. A client can
        be within its monthly quota and still be rate-limited this second.
  \item \textbf{What does HTTP 429 mean, and what header must accompany it?}
        \emph{Too Many Requests}: this caller exceeded its allowance;
        retryable later. It must carry \texttt{Retry-After} (delay-seconds
        or an HTTP-date) telling the client when to return; without it,
        clients guess, and the industry's guess is an immediate retry storm
        \cite{rfc9110}.
  \item \textbf{How does a fixed window counter work?}
        Partition time at $\lfloor t/W \rfloor$, keep one counter per key
        per window, admit while the counter is below the limit. $O(1)$
        memory, crisp reset semantics --- and the 2x boundary burst of
        Theorem~\ref{thm:ch14-fixed-window-2x}.
  \item \textbf{Explain the token bucket to a new hire.}
        A bucket holds up to $B$ tokens and refills at $r$ per second; each
        request spends one token or is rejected. $B$ is the burst you
        tolerate, $r$ is the rate you sustain, and the two knobs are
        independent --- that independence is the whole point.
  \item \textbf{Why not just store counters in each app server's memory?}
        With $R$ replicas behind a load balancer, each replica sees roughly
        $1/R$ of a client's traffic, so the effective limit becomes
        $R \times$ the configured one --- and it varies with routing
        stickiness. Shared state (a Redis cell) or a deliberately
        partitioned local budget ($\mathrm{limit}/R$) is required.
  \item \textbf{What is a concurrency limit and how is it different?}
        It bounds requests \emph{in flight}, not started per unit time. By
        Little's law, concurrency $=$ throughput $\times$ latency, so slow
        requests consume concurrency budget even when they are well within
        any rate limit. You need both knobs.
  \item \textbf{What information should a 429 response carry?}
        Status 429; \texttt{Retry-After} in delay-seconds (ceiled);
        ideally \texttt{RateLimit-Limit}, \texttt{RateLimit-Remaining: 0},
        \texttt{RateLimit-Reset}, and \texttt{RateLimit-Policy}; and a
        problem-details body (Chapter~\ref{ch:problem_details}) stating
        which limit was hit and how to get more.
  \item \textbf{Why inject a clock into a rate limiter?}
        Determinism. Time-advancing tests instead of \texttt{sleep}-based
        tests are fast, exact at boundaries, and never flaky on a loaded CI
        runner; the same injection lets you replay incident timelines
        against production code.
  \item \textbf{What is \texttt{RateLimit-Policy: "100;w=60"}?}
        The IETF draft field that states the rule itself: 100 requests per
        60-second window. Paired with \texttt{RateLimit-Remaining} on
        successes, it lets clients self-throttle before ever being
        rejected~\cite{ietf_ratelimit_headers}.
\end{enumerate}

\subsection*{Mid (Q11--Q20)}
\begin{enumerate}[label=\textbf{Q\arabic*.}, leftmargin=1.6cm, start=11]
  \item \textbf{Prove the fixed-window 2x burst flaw.}
        Fire $N$ requests at $kW - \varepsilon$ (window $k-1$ admits all)
        and $N$ more at $kW + \varepsilon$ (window $k$'s counter reset,
        admits all): $2N$ admissions inside an interval of length
        $2\varepsilon \le W$. Tight because any $W$-length interval touches
        at most two windows (Theorem~\ref{thm:ch14-fixed-window-2x}).
        Consequence: instantaneous rate is unbounded --- the $2N$ can
        arrive within milliseconds.
  \item \textbf{Token bucket versus leaky bucket --- what's actually
        different?}
        ``Leaky bucket'' is two machines. As a \emph{meter} (drain at $r$,
        overflow rejects) it is GCRA, which is provably the token bucket in
        different coordinates (Proposition~\ref{prop:ch14-gcra-equiv}): both
        admit bursts up to a tolerance and sustain $r$. As a \emph{queue}
        (FIFO drained at $r$) it is the opposite of a token bucket: bursts
        are destroyed, output is perfectly smooth, and the cost is queueing
        delay up to $Q/r$. The interview trap is answering before asking
        which machine the question means.
  \item \textbf{Derive GCRA's TAT update rule and its parameters.}
        Let $T = 1/r$ be the emission interval and $\tau$ the burst
        tolerance; keep $\mathrm{TAT}$, the theoretical arrival time of the
        next conforming request. On arrival $t_a$: reject if
        $t_a < \mathrm{TAT} - \tau$ (the request is earlier than tolerated
        jitter allows); else admit and set
        $\mathrm{TAT} \leftarrow \max(t_a, \mathrm{TAT}) + T$ --- each
        admission defers the schedule by one interval, and the $\max$
        re-anchors after idle time so credit never accrues beyond $\tau$.
        Retry-After falls out as $\mathrm{TAT} - \tau - t_a$.
  \item \textbf{When is the sliding window log worth its memory?}
        When exactness is the product: authentication attempts, payment
        authorization, license seat check-in --- anywhere ``approximately 5
        attempts'' is not a defense. Note $N$ is tiny in exactly those
        domains, so the $O(N)$ log is cheap precisely where you need it.
  \item \textbf{State and bound the sliding window counter's error.}
        Estimator $\hat{c} = c + p(1-f)$, with $p$ the previous window's
        count and $f$ the elapsed fraction. Error
        $|\hat{c} - c_{\mathrm{true}}| \le p\max(f, 1-f) < p \le N$, worst
        when last window's traffic was bunched at one end --- i.e., worst
        right after the limit mattered most
        (Proposition~\ref{prop:ch14-swc-error}).
  \item \textbf{Why is \texttt{INCR}-then-reject in Redis still wrong?}
        It is atomic but semantically broken: rejected requests consume
        increments, so an attacker (or a retrying client) inflates the
        counter and poisons the window for everyone, including itself; the
        count also stops meaning ``admitted requests,'' breaking your
        headers and metrics. Check-and-admit must be conditional inside one
        script.
  \item \textbf{Design the Redis+Lua sliding window. What commands, and why
        one script?}
        \texttt{ZREMRANGEBYSCORE} to evict, \texttt{ZCARD} to count,
        conditional \texttt{ZADD} (unique member, timestamp score),
        \texttt{PEXPIRE} for memory hygiene, all in one Lua script because
        Redis runs scripts atomically --- the evict-count-admit sequence
        can never interleave with another client's
        (Listing~\ref{lst:ch14-redis-lua}). Return remaining and a
        millisecond retry-after so the caller can emit accurate headers.
  \item \textbf{Your limiter's Redis dies at 2~a.m. Walk me through the
        options.}
        Fail-open (serve everything; abuse window until recovery; right for
        read-heavy public APIs), fail-closed (reject everything; right for
        login/payment; must exempt health checks or the load balancer
        buries you), or degraded-local (per-replica $\mathrm{limit}/R$
        fallback; bounded over-admission; the default when a local limiter
        exists). The decision belongs in a per-endpoint-class matrix made
        before the incident (Table~\ref{tab:ch14-failmatrix}).
  \item \textbf{How do you choose $B$ for a token bucket?}
        From the legitimate burst shape of your best-behaved heavy client:
        e.g., a mobile fleet that syncs 20 records on reconnect needs
        $B \ge 20$ even if its sustained rate is 2/s. Set $r$ from
        capacity, $B$ from legitimate burst, and reject the temptation to
        use $B$ to absorb \emph{illegitimate} bursts --- that is what 429
        and shedding are for.
  \item \textbf{Per-IP limiting: when does it help and when does it hurt?}
        Helps pre-authentication (login, key exchange) as a coarse outer
        net. Hurts behind NAT/CGN (one IP $=$ an enterprise: collective
        punishment) and against distributed attacks (botnets divide the
        limit by pool size). Never the only dimension --- see the case
        study (Section~\ref{sec:ch14-case}).
\end{enumerate}

\subsection*{Senior (Q21--Q30)}
\begin{enumerate}[label=\textbf{Q\arabic*.}, leftmargin=1.6cm, start=21]
  \item \textbf{Prove the token bucket's envelope and explain why it is
        tight.}
        Tokens are conserved: admissions in $[t_0, t_1]$ remove $A$ tokens,
        supply is the initial $\le B$ plus refill $r(t_1 - t_0)$, so
        $A \le B + r\Delta t$; attained by arriving as tokens accumulate.
        Tightness matters: \emph{any} limiter admitting bursts of $B$ and
        sustaining $r$ must admit at least this envelope, so the token
        bucket is the most permissive limiter with those two guarantees ---
        the exact frontier of the design space.
  \item \textbf{How would you enforce one global limit across three
        regions with no synchronous cross-region calls?}
        Regional cells with home-cell pinning: consistent-hash each key to
        a region, enforce there; accept that a client active in two regions
        gets up to $2\times$ during failover windows. If the limit is
        billing-grade, quotas must be exact: make the quota path
        single-home (one authoritative region per key) and accept its
        latency, while rate protection stays regional. State the
        consistency/latency frontier explicitly in the design doc.
  \item \textbf{Design rate limiting for a multi-tenant SaaS.}
        Layered: (1) per-IP coarse net pre-auth; (2) per-API-key token
        bucket for rate, sized by plan tier; (3) per-account monthly quota
        for billing; (4) per-endpoint cost weights so exports and batch
        jobs drain budget proportionally; (5) per-route concurrency caps on
        expensive endpoints; (6) global shedding by priority class when the
        system itself is hot. State in Redis cells, one Lua script per
        decision, consistent-hashed by tenant; noisy-neighbor isolation
        comes from (2)+(4), plan differentiation from (3), survival from
        (5)+(6).
  \item \textbf{How do you rate limit GraphQL by cost?}
        The route is constant (\texttt{POST /graphql}); the \emph{query} is
        the cost. Parse and statically score it: sum per-field weights,
        multiply by pagination first-arguments along list edges, penalize
        depth (or reject depth past a ceiling outright), charge
        introspection separately, and debit the score from a per-key token
        bucket \emph{before} execution --- ideally re-measuring after
        execution to calibrate weights. This is the query-cost analysis of
        Chapter~\ref{ch:graphql} wired into this chapter's admission
        machinery; without it, one nested query walks through a
        requests-per-minute limit shaped for REST.
  \item \textbf{What is AIMD and where does it fit in limiting?}
        Additive increase, multiplicative decrease: $L \leftarrow L + a$
        while healthy, $L \leftarrow \beta L$ ($\beta < 1$) on overload
        signals (p99 latency, 5xx, queue depth). It adapts limits to a
        capacity you cannot measure directly, and the asymmetry ---
        decrease faster than you increase --- is what makes it safe. Same
        mathematics as TCP congestion control; commonly applied to
        concurrency limits server-side and to per-tenant ceilings during
        brownouts.
  \item \textbf{Your \texttt{RateLimit-Remaining} says 3 but the next
        request is rejected. List the causes.}
        Multi-replica skew (client hit a different cell or a local
        limiter); concurrent requests from the same key (remaining is a
        point-in-time read, not a reservation); cost-weighted charging (the
        next request costs more than 1); a second, stricter dimension
        (per-endpoint or global bucket) firing; clock skew between cells;
        or a lying middleware --- which is why
        Listing~\ref{lst:ch14-middleware-tests} tests headers across
        window boundaries.
  \item \textbf{How do you size a login endpoint's defenses?}
        Per-account exact sliding log (tiny $N$, so exactness is free),
        per-IP loose outer net, global endpoint token bucket as a
        circuit-level budget, progressive delay on repeated failures, and
        velocity alerting on the global attempt rate. Fail-closed for the
        account limiter, fail-open for the IP net, and never let any of it
        429 the health check. Justify each layer against the case study's
        three errors (Section~\ref{sec:ch14-case}).
  \item \textbf{Compare load shedding and rate limiting operationally.}
        Limiting is per-identity, steady-state, contractual: clients are
        told their budget and 429 means ``you specifically, slow down.''
        Shedding is identity-agnostic, incident-state, protective: 503 +
        \texttt{Retry-After} means ``everyone, come back later.'' Shedding
        needs a priority order decided in advance (critical $>$ interactive
        $>$ batch), must shed cheaply (before parsing bodies), and must
        never conflate the two statuses or clients' backoff logic misfires.
  \item \textbf{How do you test a rate limiter?}
        Injected clock, never sleeps; prove the documented envelope
        (boundary burst, long-run rate, equivalence to a reference) as
        executable tests; replay seeded adversarial timelines; thread-storm
        the distributed path to hunt over-admission; TestClient the header
        contract across window boundaries; and test both failure policies
        with a limiter double that always raises. Listings
        \ref{lst:ch14-tests}--\ref{lst:ch14-cost-tests} are the worked
        example.
  \item \textbf{Envoy/gateway versus in-app limiting: where should limits
        live?}
        Both, for different scopes. Edge/gateway limits (Envoy local or
        global rate-limit service) stop volumetric and pre-auth abuse
        cheaply, before your app spends anything. In-app limits enforce
        business semantics the gateway cannot see: per-plan quotas,
        cost-weighted budgets, per-account attempt counts. The anti-pattern
        is either alone: gateway-only cannot price work; app-only lets
        junk traffic consume your accept queue, TLS, and parse budget.
\end{enumerate}

\subsection*{Staff (Q31--Q36)}
\begin{enumerate}[label=\textbf{Q\arabic*.}, leftmargin=1.6cm, start=31]
  \item \textbf{Design rate limiting for a platform at $10^5$ rps across
        five regions, with hard per-tenant fairness and a per-request
        limiter latency budget of 1~ms. What falls out?}
        A synchronous global cell is dead on arrival (cross-region RTT
        alone eats the budget), so: local in-process token buckets on the
        hot path, budgets sharded as $\mathrm{limit}/R$ with
        gossip/cell-sync reconciliation on a 100~ms--1~s cadence, exact
        enforcement reserved for the narrow exactness-critical surface
        (auth, payments) routed to home cells. Fairness becomes a
        \emph{bounded-error} claim: over-admission $\le$ replica count
        $\times$ sync lag in window units --- publish the bound, monitor
        against it, and reserve exactness for where the business pays for
        it.
  \item \textbf{A customer's retry storm is indistinguishable from a DDoS
        at the edge. How do you tell them apart, and how does the answer
        change what the limiter does?}
        Intent is invisible; \emph{shape} is not: customer storms carry
        valid auth, stable key sets, and requests that \emph{would} succeed
        if admitted, while attack traffic fails auth or probes randomly.
        For the customer storm, rejection must be maximally informative
        (accurate \texttt{Retry-After}, jitter guidance) because they will
        obey it; for attack traffic, cheapness dominates --- reject before
        auth, before parsing, ideally at a layer you do not pay per-request
        for. Same 429, different economics; a limiter that cannot reject
        junk cheaply is a cost amplifier.
  \item \textbf{Derive the over-admission bound for $R$ local limiters with
        sync interval $\Delta$ and window $W$, and defend or attack the
        design.}
        Each replica enforces $N/R$ locally; between syncs a key can be
        admitted $N/R$ on \emph{each} replica, so worst-case admissions per
        $W$ are $N$ plus drift: replicas observe each other's consumption
        $\Delta$ late, adding at most $R \cdot r \cdot \Delta$ where $r$ is
        the attack rate; with adversarial split-brain routing the bound
        degrades to $R \times N/R = N$ \emph{per unsynced window} plus the
        drift term. Defense: the bound is honest, tunable (shrink $\Delta$,
        over-divide $N$), and failure-safe (a dead replica's share is
        lost, not stolen). Attack: under targeted abuse the attacker routes
        around awareness, so exactness-critical surfaces must not use it.
        Both are correct; that is the point of publishing the bound.
  \item \textbf{When should limits be contractual (documented, stable,
        versioned) versus dynamic (adaptive, opaque)? What does each cost
        the ecosystem?}
        Contractual limits are an API surface: SDKs codegen backoff logic
        against them, integrators size batch jobs against them, and
        changing them is a breaking change in the Chapter~\ref{ch:versioning}
        sense --- announce, deprecate, migrate. Dynamic limits adapt to
        capacity you cannot promise, but opacity pushes integrators toward
        defensive polling and support tickets. The mature posture: contract
        the \emph{floor} (``you will always get at least $X$''), keep the
        headroom dynamic, and version the policy string itself so header
        semantics never shift under a pinned integration.
  \item \textbf{Design the migration from a broken limiting system
        (fixed window, per-IP only, no headers) to a correct one on a live
        public API.}
        Expand-and-contract (Chapter~\ref{ch:versioning}): (1) emit the new
        headers \emph{observationally} alongside the old limiter --- log
        what the new algorithms \emph{would} have done; (2) enforce the new
        limiter in shadow mode, comparing admit sets to bound behavioral
        delta; (3) tighten enforcement per tenant cohort with per-cohort
        kill switches, starting with internal keys; (4) publish the policy
        change with a deprecation window for clients whose traffic the new
        envelopes will newly reject; (5) retire the per-IP-only path last,
        after auth coverage makes it redundant. At every step, the headers
        tell clients what is coming before enforcement arrives.
  \item \textbf{Is rate limiting a security control, a reliability control,
        or a billing control? Argue all three, then prioritize.}
        Reliability: it is the cheap-rejection valve that keeps offered
        load below the capacity knee --- the primary justification.
        Security: it raises the cost of credential stuffing, scraping, and
        inventory abuse, but it is a \emph{friction} control, not a
        barrier --- authentication, authorization, and detection remain the
        barriers (Chapter~\ref{ch:api_security}). Billing: quotas encode
        the business model, and metering-grade exactness there is a
        different (harder, single-home) engineering problem than
        protection-grade approximation. Prioritize reliability ---
        a down API has no security or revenue --- but design the dimensions
        so security and billing ride the same machinery with their own
        exactness budgets.
\end{enumerate}
%% ----------------------------------------------------------------------------
\section{External Bibliography \& Further Reading}
\label{sec:ch14-bibliography}

\begin{itemize}
  \item IETF HTTPAPI Working Group, \emph{RateLimit Header Fields for HTTP}
        (Internet-Draft) --- the standardizing vocabulary for
        \texttt{RateLimit-Limit/-Remaining/-Reset} and
        \texttt{RateLimit-Policy}; read the draft's deployment notes for
        the epoch-versus-delay-seconds reset
        history~\cite{ietf_ratelimit_headers}.
  \item RFC~9110, \emph{HTTP Semantics} --- \texttt{Retry-After} (both
        forms), 429's relationship to the 4xx family, and 503 semantics for
        overload~\cite{rfc9110}.
  \item RFC~6585, \emph{Additional HTTP Status Codes} --- where 429 is
        actually registered (uncited in text; short and worth reading in
        full).
  \item M.~Nygard, \emph{Release It!}, 2nd ed. --- capacity knees, the
        stability patterns (bulkheads, circuit breakers) that rate limiting
        complements, and why unbounded queues fail
        slowly~\cite{nygard2018releaseit}.
  \item M.~Kleppmann, \emph{Designing Data-Intensive Applications} --- the
        shared-state and replication background behind
        Section~\ref{subsec:ch14-distributed}'s atomicity and cell
        arguments~\cite{kleppmann2017ddia}.
  \item Cloudflare, \emph{Counting things: a lot of different things}
        (engineering blog) --- the sliding-window-counter error analysis on
        real traffic that Proposition~\ref{prop:ch14-swc-error} formalizes.
  \item Stripe, \emph{Rate limiters} (engineering blog) --- production
        token-bucket fleet design, including the load-shedder priority
        classes that inspired Section~\ref{subsec:ch14-multidim}.
  \item Kong, \emph{How to Design a Scalable Rate Limiting Algorithm}
        (engineering blog) --- the gateway-level view of the same
        algorithm zoo, including the cluster-versus-local trade-off.
  \item NGINX, \emph{Rate Limiting with NGINX and NGINX Plus} --- the
        \texttt{limit\_req} documentation where GCRA hides in plain sight
        (the \texttt{rate} and \texttt{burst} parameters map to $T$ and
        $\tau$).
  \item Redis documentation, \emph{Rate limiter patterns} --- the canonical
        \texttt{INCR}/\texttt{EXPIRE} and sorted-set recipes, including the
        atomicity caveats our Lua script resolves.
  \item ITU-T I.371 / ATM Forum TM~4.0 --- the original Generic Cell Rate
        Algorithm specification; historically illuminating for how little
        the problem has changed.
  \item FastAPI documentation --- middleware and \texttt{TestClient} used
        throughout Listings 14.6--14.7~\cite{fastapidocs}.
\end{itemize}

%% ----------------------------------------------------------------------------
\section{Capstone Projects}
\label{sec:ch14-capstones}

\subsection*{Capstone 14.1 --- Prove the Algorithms \hfill [junior]}
\textbf{Objective:} Reproduce this chapter's entire proof sheet from a
standing start, then break it on purpose.

\textbf{Inputs/Datasets:} Listings 14.1--14.2; no datasets.

\textbf{Steps:} (1)~Reimplement the five limiters from their mathematical
definitions only (close the book; keep the \texttt{Decision} contract).
(2)~Port the 14 tests of Listing~\ref{lst:ch14-tests} to your
implementation unmodified. (3)~Introduce one deliberate bug per algorithm
(e.g., forget the $\max$ in GCRA's TAT update) and record which test
catches each. (4)~Write one new test: the fixed window's 2x burst with a
\emph{non-integer} window (e.g., $W = 1.5$).

\textbf{Acceptance Criteria:}
\begin{itemize}
  \item[$\square$] All 14 ported tests pass against your reimplementation.
  \item[$\square$] Every seeded bug is caught by at least one existing
        test; you can name the mapping.
  \item[$\square$] The new non-integer-window burst test proves 20
        admissions in a sub-window interval.
\end{itemize}

\textbf{Stretch Goals:} Property-test the sliding log with random seeded
timelines, asserting Proposition~\ref{prop:ch14-swl-exact}'s invariant on
every run.

\subsection*{Capstone 14.2 --- Header-Complete Middleware \hfill [mid]}
\textbf{Objective:} Extend Listing~\ref{lst:ch14-middleware} into middleware
you would publish as a platform library.

\textbf{Inputs/Datasets:} Listings 14.1, 14.6, 14.7; the Northwind Robotics
endpoint set (\texttt{/telemetry}, \texttt{/fleet/status},
\texttt{/reports/export}, \texttt{/health}).

\textbf{Steps:} (1)~Swap the in-process limiter for the Redis limiter of
Listing~\ref{lst:ch14-redis} behind the \texttt{KeyedLimiter} protocol
(add a \texttt{reset\_after} implementation). (2)~Emit an RFC~9457-style
problem body (Chapter~\ref{ch:problem_details}) on 429, including the
policy and a documentation link. (3)~Add an \texttt{X-NR-Limiter-Degraded:
true} signal for fail-open episodes (headers omitted, signal present).
(4)~Test: headers across a window boundary, problem body shape, degraded
signal, and 401 versus 429 separation.

\textbf{Acceptance Criteria:}
\begin{itemize}
  \item[$\square$] All Listing~\ref{lst:ch14-middleware-tests} tests pass
        unmodified against the Redis-backed wiring.
  \item[$\square$] 429 bodies validate against a declared problem-details
        schema.
  \item[$\square$] Fail-open episodes emit the degraded signal and no
        \texttt{RateLimit-*} fields; fail-closed returns 503 +
        \texttt{Retry-After}.
  \item[$\square$] Full suite is deterministic and offline.
\end{itemize}

\textbf{Stretch Goals:} Add per-plan policies (\texttt{"100;w=60"} vs.\
\texttt{"1000;w=60"}) selected by API key tier.

\subsection*{Capstone 14.3 --- Cell Architecture Simulator \hfill [mid]}
\textbf{Objective:} Quantify, rather than assert, the over-admission of
cell-based limiting.

\textbf{Inputs/Datasets:} Listing~\ref{lst:ch14-sim}'s harness pattern; a
seeded multi-client timeline generator you write (10 keys, Poisson-ish
arrivals, one adversarial key).

\textbf{Steps:} (1)~Implement three enforcement topologies over the same
timeline: one global limiter (exact), $R=4$ local limiters at
$\mathrm{limit}/R$ with no sync, and $R=4$ local limiters with a sync every
$\Delta$ seconds (reconciliation subtracts observed remote consumption).
(2)~Sweep $\Delta \in \{0.1, 0.5, 1.0, 5.0\}$ s and record worst-case
per-window admissions for the adversarial key. (3)~Produce a pgfplots-ready
table of over-admission versus $\Delta$. (4)~Verify the
Q33-style bound: admissions $\le N + R \cdot r \cdot \Delta$.

\textbf{Acceptance Criteria:}
\begin{itemize}
  \item[$\square$] Global topology never exceeds the limit; no-sync
        topology's overshoot matches $R \times$ prediction.
  \item[$\square$] Sync topology's measured overshoot respects the
        published bound at every $\Delta$.
  \item[$\square$] All results reproducible from one seed.
\end{itemize}

\textbf{Stretch Goals:} Add a cell-failure event mid-timeline and show the
consistent-hash remap confines counter loss to $1/C$ of keys.

\subsection*{Capstone 14.4 --- Adaptive Limits with AIMD \hfill [senior]}
\textbf{Objective:} Build and tune an AIMD-adaptive per-tenant limiter that
tracks a moving capacity.

\textbf{Inputs/Datasets:} Listing~\ref{lst:ch14-limiters}; a synthetic
``capacity signal'' generator you author: service capacity oscillates
between 150 and 450 rps on a seeded schedule, and overload is reported when
admitted rate exceeds current capacity.

\textbf{Steps:} (1)~Implement the controller: per-second evaluation,
$L \leftarrow L + a$ when healthy, $L \leftarrow \beta L$ on overload,
floored at a contractual minimum. (2)~Sweep $(a, \beta)$ over a small grid;
for each pair record time-to-first-contact, steady-state oscillation
amplitude, and total rejected load. (3)~Identify the knee: the fastest
$a$ whose $\beta$-paired oscillation never exceeds capacity by more than
10\%. (4)~Write the tuning memo: why multiplicative decrease is the safety
parameter and additive increase merely the efficiency parameter.

\textbf{Acceptance Criteria:}
\begin{itemize}
  \item[$\square$] Controller never lets admitted rate exceed capacity by
        $>10\%$ for more than one evaluation interval, for the tuned
        $(a, \beta)$.
  \item[$\square$] Floor is never violated even under sustained overload.
  \item[$\square$] Grid results and the knee are reproducible from the
        seed; memo quantifies the asymmetry argument.
\end{itemize}

\textbf{Stretch Goals:} Add per-tenant fairness: divide recovered headroom
proportional to recent demand rather than equally.

\subsection*{Capstone 14.5 --- Rate-Limiting Architecture for Northwind at
100k rps \hfill [staff]}
\textbf{Objective:} Produce the complete design document for rate limiting,
quotas, and shedding across Northwind Robotics' public API surface at
$10^5$ rps in five regions.

\textbf{Inputs/Datasets:} This chapter; a synthetic endpoint inventory you
author (40 endpoints across read/write/export/admin classes with measured
cost weights); a synthetic tenant registry (3 tiers, 2{,}000 tenants, one
known-abusive actor).

\textbf{Steps:} (1)~Classify every endpoint into a protection tier
(exactness-required, approximate-fair, shedding-only) and defend each
classification. (2)~Choose the topology per tier (global cell / regional /
local-plus-sync) with latency budgets; justify the 1~ms hot-path budget's
consequences. (3)~Specify the header contract, quota endpoints, and the
policy-versioning story. (4)~Write the failure runbook: limiter-cell
outage, abusive-tenant drill, capacity-brownout (AIMD kick-in), including
the fail-open/fail-closed assignment per tier. (5)~Model cost: Redis cell
sizing from key cardinality and $N$, and the egress cost of rejecting
cheaply versus serving slowly. (6)~Red-team your own design with the
Section~\ref{sec:ch14-case} playbook.

\textbf{Acceptance Criteria:}
\begin{itemize}
  \item[$\square$] Every endpoint has a tier, a topology, a failure policy,
        and a numeric limit with a derivation (not a round number from
        nowhere).
  \item[$\square$] Over-admission bounds are stated and monitored per
        approximate tier.
  \item[$\square$] Runbook covers the three named incidents with detection
        signals and rollback steps.
  \item[$\square$] Cost model shows limiter infrastructure $< 2\%$ of
        serving cost at $10^5$ rps, or explains why not.
  \item[$\square$] Red-team section names at least two residual weaknesses
        honestly.
\end{itemize}

\textbf{Stretch Goals:} Specify the GraphQL cost-scoring integration
(Chapter~\ref{ch:graphql}) and the gRPC interceptor analogue
(Chapter~\ref{ch:grpc}) so all three protocol surfaces share one budget.

%% ----------------------------------------------------------------------------
\section{Python Dependencies Appendix}
\label{sec:ch14-deps}

All code runs offline after a one-time install. From PowerShell:

\begin{minted}{text}
python -m venv .venv
.\.venv\Scripts\Activate.ps1
pip install -r requirements.txt
\end{minted}

\noindent\textbf{requirements.txt:}
\begin{minted}{text}
fakeredis>=2.23
fastapi>=0.110
httpx>=0.27
pytest>=8
\end{minted}

For real deployments (never needed by this chapter's listings): add
\texttt{redis>=5} and swap \texttt{fakeredis.FakeRedis()} for
\texttt{redis.Redis(...)} --- the Lua script, calling convention, and
\texttt{Decision} contract are identical.

\subsection{fakeredis ($\geq$ 2.23)}
\textbf{Description:} In-process emulation of the Redis server, including
Lua script execution. \textbf{Why included:} lets the distributed limiter
(Listings 14.4--14.5) and its thread-storm atomicity test run with zero
infrastructure --- the same code paths, the same script, no server, no
containers. \textbf{Alternatives:} a real Redis in Docker (faithful,
including wall-clock expiry timing, at the price of infrastructure);
\texttt{testcontainers}-based fixtures (CI-friendly, still heavyweight);
mocking \texttt{eval} at the client boundary (fast but verifies nothing
about the script). \textbf{Installation:}
\texttt{pip install "fakeredis>=2.23"} (Lua support requires the bundled
\texttt{lupa} dependency; the pinned distribution includes it).

\subsection{fastapi ($\geq$ 0.110)}
\textbf{Description:} The ASGI web framework hosting the middleware
demonstrations. \textbf{Why included:} routing, dependency injection, and
the \texttt{TestClient} that makes Listing~\ref{lst:ch14-middleware}'s
header assertions runnable offline~\cite{fastapidocs}.
\textbf{Alternatives:} Starlette directly (FastAPI's base; the middleware
is raw ASGI and would work unchanged); Flask (WSGI --- the ASGI middleware
pattern would need reworking); a gateway doing the limiting instead (Envoy
local rate limit, Kong) for the edge layer. \textbf{Installation:}
\texttt{pip install "fastapi>=0.110"}.

\subsection{httpx ($\geq$ 0.27)}
\textbf{Description:} Sync/async HTTP client. \textbf{Why included:} it is
the transport under FastAPI's \texttt{TestClient}; required for the
middleware tests to execute without a server. \textbf{Alternatives:}
\texttt{requests} (synchronous, not ASGI-native); running a live
\texttt{uvicorn} and a real socket (faithful but no longer offline-pure).
\textbf{Installation:} \texttt{pip install "httpx>=0.27"}.

\subsection{pytest ($\geq$ 8)}
\textbf{Description:} Test framework. \textbf{Why included:} runs all 35
chapter tests, including the executable proofs and the concurrency storm.
\textbf{Alternatives:} \texttt{unittest} (stdlib; works, but fixtures and
\texttt{approx} would need hand-rolling); \texttt{nose2} (unmaintained
momentum). \textbf{Installation:} \texttt{pip install "pytest>=8"}.

\subsection*{Ecosystem alternatives worth knowing}
For application-level limiting beyond this chapter's teaching
implementations: \texttt{limits} (mature strategy library: fixed/sliding
algorithms over Redis/Memcached/in-memory backends), \texttt{slowapi}
(FastAPI/Starlette integration built on \texttt{limits}),
\texttt{django-ratelimit} (Django decorator), and Envoy's
\emph{local rate limit} filter / global rate-limit service for the edge
tier. Evaluate them against Section~\ref{sec:ch14-mechanics}: which
algorithm, which atomicity story, which failure policy, which headers.

%% ----------------------------------------------------------------------------
\section{Chapter Summary \& Key Takeaways}
\label{sec:ch14-summary}

\begin{itemize}
  \item Five tools, five questions: rate limiting (how fast may this caller
        start), throttling (shape the flow), quotas (how much in total),
        concurrency limits (how many in flight), load shedding (which work
        do we drop). Deploying one and claiming the others is the canonical
        maturity failure.
  \item The fixed window's 2x boundary burst is a theorem with a tight
        bound, not folklore (Theorem~\ref{thm:ch14-fixed-window-2x}); if
        the limit protects something with a hard per-second ceiling, the
        fixed window does not enforce it.
  \item The sliding log is exact at $O(N)$ memory; the sliding counter
        buys $O(1)$ at error $<$ previous-window count
        (Proposition~\ref{prop:ch14-swc-error}), worst exactly after
        saturation; the token bucket owns the tightest burst-plus-rate
        envelope $A \le B + r\Delta t$
        (Proposition~\ref{prop:ch14-tb-envelope}); GCRA is the token bucket
        in one timestamp (Proposition~\ref{prop:ch14-gcra-equiv}).
  \item Distribution moves the problem from algorithms to state: atomic
        check-and-admit via one Lua script in one cell; consistent hashing
        to bound remap on cell loss; local-plus-sync trades exactness for
        latency with a publishable over-admission bound.
  \item The failure policy is a business decision made in advance, per
        endpoint class: fail-open, fail-closed, or degraded-local
        (Table~\ref{tab:ch14-failmatrix}); a fail-closed limiter without a
        health-check exemption is a self-inflicted outage.
  \item Headers are the contract: 429 + \texttt{Retry-After}
        (delay-seconds, ceiled), \texttt{RateLimit-Limit/-Remaining/-Reset/-Policy}
        on successes too, 503 for shedding, quota endpoints for budgets.
        Omit what you cannot compute truthfully; never fake it
        \cite{ietf_ratelimit_headers,rfc9110}.
  \item Multi-dimensional is the default at maturity: per-key for fairness,
        per-IP as a coarse pre-auth net (never the only net ---
        Section~\ref{sec:ch14-case}), per-endpoint and cost-weighted for
        expensive routes, AIMD for adapting to unmeasurable capacity,
        priority-aware shedding for survival.
  \item Test limiters with injected clocks, seeded timelines, thread
        storms, and a limiter double that always raises --- then both
        failure paths are tested before the 3~a.m.\ page tests them for you.
\end{itemize}

%% ----------------------------------------------------------------------------
\section{Quick-Reference Card}
\label{sec:ch14-quickref}

\subsection*{Algorithm Matrix (one line each)}
Fixed window --- $O(1)$, crisp reset, $2N$ boundary burst $\bullet$
sliding log --- exact, $O(N)$ memory, Retry-After = oldest expiry $\bullet$
sliding counter --- $O(1)$, error $< p$, smooths boundaries $\bullet$
token bucket --- burst $B$, sustained $r$, envelope $B + r\Delta t$
$\bullet$ GCRA --- one timestamp, reject iff $t_a < \mathrm{TAT} - \tau$,
$\mathrm{TAT} \leftarrow \max(t_a, \mathrm{TAT}) + T$, $B = 1 + \tau/T$.

\subsection*{Header Reference}
429 \emph{Too Many Requests} --- caller over budget; always with
\texttt{Retry-After: <delay-seconds>} (ceiled; HTTP-date is legal, worse)
$\bullet$ \texttt{RateLimit-Limit} --- window quota $\bullet$
\texttt{RateLimit-Remaining} --- budget left after this response $\bullet$
\texttt{RateLimit-Reset} --- delay-seconds to full reset $\bullet$
\texttt{RateLimit-Policy} --- the rule, e.g.\ \texttt{"100;w=60"} $\bullet$
503 + \texttt{Retry-After} --- system shedding, not the caller's fault
$\bullet$ 401/403 --- no key, no bucket: never 429 an anonymous caller.

\subsection*{Decision Tree (compressed)}
Shared limit across replicas? $\to$ Redis cell, one Lua script; hash keys
to cells $\bullet$ Exactness required? $\to$ sliding log $\bullet$ Bursts
legitimate? $\to$ token bucket / GCRA, $B$ from client burst shape $\bullet$
Huge smooth keyspace? $\to$ sliding counter $\bullet$ Crisp headers, rough
protection? $\to$ fixed window, short $W$ $\bullet$ Limiter down? $\to$
per-class policy: fail-open (reads), fail-closed (auth/payment),
degraded-local (default) $\bullet$ Expensive route? $\to$ cost weight +
concurrency cap + consider 202-async.

\section{Standardized Evidence Addendum}
\subsection{Benchmark Snapshot Template}
\begin{table}[h]
  \centering
  \small
  \begin{tabularx}{\textwidth}{l c c c c}
    \toprule
    \textbf{Config} & \textbf{p50 latency} & \textbf{p99 latency} & \textbf{Error rate} & \textbf{Throughput} \\
    \midrule
    Baseline & -- & -- & -- & 1.00x \\
    Candidate A & -- & -- & -- & -- \\
    Candidate B & -- & -- & -- & -- \\
    \bottomrule
  \end{tabularx}
  \caption{Minimum benchmark table to keep per chapter.}
\end{table}

\subsection{Request Trace Template}
\begin{itemize}
  \item \textbf{Request:} one representative API call for this chapter's topic.
  \item \textbf{Before:} behavior (headers, payload, latency, failure mode) with the naive approach.
  \item \textbf{After:} behavior with the technique in this chapter.
  \item \textbf{Result:} what changed in correctness, resilience, latency, and operability.
\end{itemize}

\subsection{Cost and Latency Budget Snapshot}
\begin{table}[h]
  \centering
  \small
  \begin{tabularx}{\textwidth}{l c c}
    \toprule
    \textbf{Stage} & \textbf{Latency budget} & \textbf{Cost driver} \\
    \midrule
    TLS / connection setup & -- & handshake round trips, keep-alive hit rate \\
    Request validation & -- & schema complexity, CPU per request \\
    Business logic and I/O & -- & downstream fan-out, query cost \\
    Serialization and egress & -- & payload size, compression ratio \\
    \bottomrule
  \end{tabularx}
  \caption{Operational budget template for this chapter's technique.}
\end{table}

\section{Configuration Preset Template}
\begin{minted}{text}
# api_policy.yaml
policy_version: "v1.0.0"
component: "<chapter_topic>"
timeout_ms: 0
retry_budget: 0
fallback_strategy: "fail_fast"
rollback_trigger: "error_budget_burn_gt_threshold"
\end{minted}

\section{CI Quality Gate Specification}
\begin{itemize}
  \item contract tests green against the pinned OpenAPI/AsyncAPI/Protobuf spec,
  \item no breaking schema changes without a version bump,
  \item error responses conform to RFC 9457 problem-details shape,
  \item benchmark deltas within approved regression bounds (latency and throughput),
  \item policy version and rollback plan present in release metadata.
\end{itemize}

\section{Incident Runbook Template}
\begin{enumerate}
  \item Detect regression from SLO burn-rate alerts or consumer reports.
  \item Scope impacted endpoints, versions, and consumer segments.
  \item Contain impact (rollback, feature flag, rate-limit tightening, or traffic shift).
  \item Diagnose with traces, structured logs, and contract diffs.
  \item Recover with validated fix and verified rollout procedure.
  \item Prevent recurrence via new regression tests and CI gates.
\end{enumerate}

\section{Cross-Chapter Decision Card}
\begin{table}[h]
  \centering
  \small
  \begin{tabularx}{\textwidth}{l X X}
    \toprule
    \textbf{Dimension} & \textbf{Choose this approach when} & \textbf{Prefer an alternative when} \\
    \midrule
    Use case & -- & -- \\
    Latency budget & -- & -- \\
    Operational cost & -- & -- \\
    Team maturity & -- & -- \\
    \bottomrule
  \end{tabularx}
  \caption{Decision card to complete per chapter.}
\end{table}

\section{ROI Model and Build-vs-Buy}
\begin{itemize}
  \item \textbf{Build when:} the capability is differentiating, compliance forbids vendors, or scale makes vendor pricing irrational.
  \item \textbf{Buy when:} the capability is commodity (auth, gateways, observability), time-to-market dominates, or staffing is thin.
  \item \textbf{Hidden costs to model:} on-call load, upgrade cadence, expertise ramp, data egress fees.
\end{itemize}

\section{Disaster Recovery Notes (RPO/RTO)}
\begin{itemize}
  \item Define recovery point objective (RPO): maximum acceptable data loss window.
  \item Define recovery time objective (RTO): maximum acceptable restoration time.
  \item Test restores regularly; an untested backup is not a backup.
\end{itemize}

